\documentclass{article}
\usepackage[utf8]{inputenc}
\usepackage{graphicx}
\usepackage{amsmath}
\usepackage{amssymb}
\usepackage{url}
\usepackage{tabularx}
\usepackage[colorlinks=true,allcolors=blue]{hyperref}
\usepackage{caption}
\usepackage{pdfpages}
\usepackage{multicol}
\usepackage[ngerman]{babel}
\usepackage{multicol}
\usepackage{geometry}	% [textwidth=17cm,textheight = 24.5cm]
% Um die Dokumentränder zu sehen: [showframe] vor {geometry} einfügen
\usepackage{bbold}
\usepackage{makecell}
\usepackage{float}
\usepackage{braket}
\usepackage{cancel}
\usepackage{braket}
\usepackage{mathtools}
\usepackage{xltabular}
\usepackage{float}
\usepackage{MnSymbol}
\usepackage{multirow}
\usepackage{lscape}

\renewcommand{\cellalign}{tl}

\setlength{\columnsep}{1cm}
\setlength\parindent{0pt}

\newcommand{\diff}{\mathrm{d}}
\newcommand{\konst}{\mathrm{konst.}}
\newcommand{\lr}[1]{\left(#1\right)}
\newcommand{\lreck}[1]{\left[#1\right]}
\newcommand{\lrschweif}[1]{\left\{#1\right\}}
\newcommand{\dkon}[3]{\left(\frac{\partial #1}{\partial #2}\right)_{#3}}

\definecolor{gruen}{HTML}{038133}

\newcommand{\commenttext}[1]{\textcolor{gruen}{\text{#1}}}
\newcommand{\comment}[1]{{\color{gruen} {#1}}}

\allowdisplaybreaks

\begin{document}

\begin{center}
	\Large{\textbf{Laser Physics}}\\
	\small{613348}\\
	\small{\today}
\end{center}

\tableofcontents


\newpage


\section{Overview}
A Laser (\textit{Light Amplification by Stimulated Emission of Radiation}) is an optical instrument which has the properties:
	\begin{itemize}
		\item \textbf{monochromaticity}


		\item \textbf{time coherence}\\
		Measure of the correlation of a wave with itself at a later time considering a fixed point in time.\\
		Laser light is temporaly coherent, meaning it has a large coherence time/length.


		\item \textbf{spatial coherence}\\
		Measure of the correlation of a wave at point $\vec{r}_1$ and wave at point $\vec{r}_2$ considering the same point in time.\\
		Laser light is spatialy coherent, meaning it has a large coherence plane.


		\item \textbf{high intensity} \textcolor{gruen}{(can be reached)}


		\item \textbf{short pulse duration} \textcolor{gruen}{(can be reached)}
	\end{itemize}


\begin{figure}[h]
	\centering
	\fbox{
	\includegraphics[width=0.6\textwidth]{Grafiken/laser overview.pdf}
	}
	\caption{Simplified depiction of a laser}
\end{figure}


A laser consists of a
	\begin{itemize}
		\item \textbf{gain medium}\\
		Converts energy to light with a intensity dependend gain coefficient
			\begin{equation}
				g(I) = \frac{g_0}{1 + \frac{I}{I_\text{sat}}}
			\end{equation}


		\item \textbf{optical resonator}\\
		Recycles the light and choses one specific frequency


		\item \textbf{pump}\\
		Acts as an energy source of the system
	\end{itemize}


\newpage


The basic light amplification in a laser can be described by turning the change in intensity after one roundtrip of the light inside the resonator
			\begin{align}
				I(t+t_r) &= R_1R_2G(I) \cdot I,
			\end{align}

where $G(I) = 1 + g(I)$ into a differential equation
			\begin{align}
				\frac{\diff I}{\diff t} &= \lreck{ \frac{ \lr{ 1 + \frac{g_0}{1 + I/I_\text{sat}} } R_1 R_2 - 1 }{\tau_\text{rt}} } I(t),
			\end{align}

which has to be solved numerically.\\
\\
From figure \ref{overview intensity} one can see, that after a certain time, the intensity converges to a \textbf{steady state intensity}
			\begin{equation}
				I_\text{ss} = \frac{ (g_0 + 1)R_1 R_2 - 1 }{ 1 - R_1 R_2 } I_\text{sat}.
			\end{equation}
		
The laser output is then described by
	\begin{equation}
		I_\text{out} = (1 - R_2) I_\text{ss}.
	\end{equation}

\begin{figure}[h]
	\centering
	\label{overview intensity}
	\includegraphics[width=0.6\textwidth]{Grafiken/steady state intensity.pdf}
	\caption{Qualitative plot of the light intensity over time}
\end{figure}


\newpage


{\color{violet}
\hrule
	\fbox{Derivation}\\
	Assume that the lights' intensity has increased/decreased by the factor $G(I) R_1 R_2$ after one roundtrip between the mirrors. Here, $G(I) = 1 + \frac{g_0}{1 + I/I_\text{sat}}$ describes the intensity dependent effect of the gain medium and $R_1, R_2$ the reflectivities of the mirrors.\\
	Thus, the intensity after the time of a roundtrip $t_\text{rt}$ is
		\begin{align*}
			I(t + t_\text{rt}) = \lr{ 1 + \frac{g_0}{1 + I/I_\text{sat}} } R_1 R_2 I(t).
		\end{align*}
	
	Assuming negligeably small round trip times, the left hand side can be taylor expanded, which yields the following differential equation
		\begin{align*}
			I(t + t_\text{rt}) \approx I(t) + \frac{\diff I}{\diff t} t_\text{rt} &= \lr{ 1 + \frac{g_0}{1 + I/I_\text{sat}} } R_1 R_2 I(t) \\
			\Leftrightarrow \frac{\diff I}{\diff t} &= \lreck{ \frac{ \lr{ 1 + \frac{g_0}{1 + I/I_\text{sat}} } R_1 R_2 - 1 }{ t_\text{rt} } } I(t), \qquad \blacksquare
		\end{align*}
	
	which can be solved numerically.\\
	\\
	The steady state intensity is convergence value of the intensity for $t\to\infty$ and can be found by inserting $\frac{\diff I}{\diff t}$ into the previously found differential equation:
		\begin{align*}
			0 &= \lreck{ \frac{ \lr{ 1 + \frac{g_0}{1 + I(t)/I_\text{sat}} } R_1 R_2 - 1 }{ t_\text{rt} } } I(t) \\
			\stackrel{\comment{\cdot t_\text{rt} (1 + I/I_\text{sat})}}{\Leftrightarrow} 0 &= ((1 + I/I_\text{sat}) + g_0) R_1 R_2 I - I (1 + I/I_\text{sat}) = \\
			&= R_1 R_2 I + \frac{R_1 R_2}{I_\text{sat}} I + g_0 R_1 R_2 I - I - I^2 \frac{1}{I_\text{sat}} = \\
			&= I^2 \lr{ \frac{R_1R_2}{I_\text{sat}} - \frac{1}{I_\text{sat}} } + I \lr{ R_1 R_2 + g_0 R_1 R_2 - 1 } = \\
			&= I^2 + I \cdot \frac{(g_0 + 1)R_1 R_2 - 1}{R_1 R_2 - 1} I_\text{sat} \\
			\Rightarrow I_\pm &= - \frac{ \frac{(g_0 + 1)R_1 R_2 - 1}{R_1 R_2 - 1} I_\text{sat} }{2} \pm \sqrt{ \lr{ \frac{ \frac{(g_0 + 1)R_1 R_2 - 1}{R_1 R_2 - 1} I_\text{sat} }{2} }^2 } \\
			\Rightarrow I_\text{ss} &\equiv I_- = - \frac{(g_0 + 1)R_1 R_2 - 1}{R_1 R_2 - 1} I_\text{sat} = \frac{(g_0 + 1)R_1 R_2 - 1}{1 - R_1 R_2} I_\text{sat} \qquad \blacksquare
		\end{align*}
\hrule
}


\newpage


\section{Light-Matter Interaction}
\subsection{Classical Description}
The classical description describes
	\begin{enumerate}
		\item light as an electromagnetic wave
		\item matter as polarisable particles\\
		Precisely: The \textbf{driven harmonic oscillator} to model an electron bound to its nucleus
	\end{enumerate}


\begin{table}[h]
	\centering
	\caption{Overview results classical description}
	\begin{tabularx}{\textwidth}{X|X}
		\hline\hline
			Successes & Failures \\
		\hline
			- describes the amplification of the E-field in a doped medium & - doesn't describe a saturation intensity \\
		\hline
			& - needs phenomenological damping parameter \\
		\hline\hline	
	\end{tabularx}
\end{table}


Auxiliary assumptions are, that the medium is
	\begin{itemize}
		\item \textbf{linear}
		
		
		\item \textbf{isotropic}
		
		
		\item \textbf{non-magnetic}
		
		
		\item \textbf{no free charges/currents}
	\end{itemize}

as well that the E-field is
	\begin{itemize}
		\item \textbf{slowly varying}
		
		
		\item \textbf{paraxial}\\
		The field is close and almost parallel to the optical axis
		
		
		\item \textbf{monochromatic}
	\end{itemize}

The wave equation for the E-field under these assumptions is called the \textbf{inhomogeneous, paraxial wave equation}
	\begin{align*}
		\Delta_t \vec{E}_0(\vec{r},t) + 2 i n k_0 \lr{ \partial_z + \frac{\alpha}{2} + \frac{\kappa}{2} + \frac{1}{c_g}\partial_t } \vec{E}_0(\vec{r},t) &= \exp\lrschweif{ i(\omega t - n k_0 t) } \mu_0 \partial_t^2 \vec{P}_A
	\end{align*}

with the \textbf{absorption coefficient} $\alpha := k \cdot \Im\chi_M$, \textbf{gain coefficient} $g := k \cdot \Im\chi_A$, and a loss factor for additional losses $\kappa$ \textcolor{gruen}{(added phenomenologically)}.\\
\\
The solution for a \textbf{monochromatic wave} takes the form
	\begin{equation}
		\vec{E}(z,t) \propto \exp \lrschweif{ - \frac{1}{2}\lr{ \alpha + \kappa - g } \cdot z }.
	\end{equation}


\newpage


{\color{violet}
\hrule
	\fbox{Derivation}\\
	Strategy:\\
	1) Use the driven harmonic oscillator to model the motion of an electron around its nucleus.\\
	2) Formulate the polarisation of the medium.\\
	3) Then, derive the wave equation for the electric field in the medium.\\
	4) Use a lot of approximation as well as phenomenological terms to arrive at the defining equation of the E-field.\\
	\\
	The equation of motion for an electron which moves in the atom's harmonic potential, experiences friction as well as the force due to an external electric field is
		\begin{align*}
			\frac{\diff^2\vec{x}}{\diff t^2} + \frac{1}{\tau} \frac{\diff \vec{x}}{\diff t} + \omega_R^2 \vec{x} &= \frac{e}{m} \vec{E}_0 e^{-i\omega t}.
		\end{align*}
	
	The solution to this differential equation is known from classical mechanics
		\begin{align*}
			\vec{x}(t) &= \frac{e/m}{ (\omega_R^2 - \omega^2) - i\frac{\omega}{\tau} } \vec{E}_0 e^{-i\omega t}.
		\end{align*}

	Thus, the induced dipole moment of the atom is
		\begin{align*}
			\vec{p} &= e\vec{x} = \frac{e^2/m}{ (\omega_R^2 - \omega^2) - i\frac{\omega}{\tau} } \vec{E}_0 e^{-i\omega t}.
		\end{align*}

	Since the medium consists of multiple atoms, the macroscopic polarisation can be formulated using the atomic density $n_0$
		\begin{align*}
			\vec{P} &= n_0 \cdot \vec{p} = n_0 \cdot e\vec{x} = \frac{e^2/m}{ (\omega_R^2 - \omega^2) - i\frac{\omega}{\tau} } \vec{E}_0 e^{-i\omega t} \equiv \\
			&\equiv \varepsilon_0 \chi \vec{E},
		\end{align*}

	where in the last step, the assumption of a linear medium was used and the susceptibility defined as
		\begin{align*}
			\chi &:= \frac{n_0 e^2}{\varepsilon_0 m} \frac{1}{ (\omega_R^2 - \omega^2) - i\frac{\omega}{\tau} } = \\
			&\stackrel{\comment{1}}{=} \frac{n_0 e^2}{\varepsilon_0 m} \frac{1}{ (\omega_R^2 - \omega^2) - i\frac{\omega}{\tau} } \cdot \frac{ (\omega_R^2 - \omega^2) + i\frac{\omega}{\tau} }{ (\omega_R^2 - \omega^2) + i\frac{\omega}{\tau} } = \\
			&= \frac{n_0 e^2}{\varepsilon_0 m} \frac{ (\omega_R^2 - \omega^2) + i\frac{\omega}{\tau} }{ (\omega_R^2 - \omega^2)^2 + (\frac{\omega}{\tau})^2 } = \\
			&= \frac{n_0 e^2}{\varepsilon_0 m} \frac{ (\omega_R^2 - \omega^2)}{ (\omega_R^2 - \omega^2)^2 + (\frac{\omega}{\tau})^2 } + i\frac{n_0 e^2}{\varepsilon_0 m} \frac{\frac{\omega}{\tau} }{ (\omega_R^2 - \omega^2)^2 + (\frac{\omega}{\tau})^2 } \equiv \\
			&\equiv \chi_R + i \chi_I
		\end{align*}


\newpage


	Now, derive the wave equation for the linear, isotropic, free-charge and free-current-free medium.\\
	The Maxwell equations under these assumptions simplify to
		\begin{align*}
			\vec{\nabla} \vec{D} &= 0 \\
			\vec{\nabla} \times \vec{E} &= -\partial_t \vec{B} \\
			\vec{\nabla} \vec{B} &= 0 \\
			\vec{\nabla} \times \vec{H} &= \partial_t \vec{D}
		\end{align*}

	First express $\vec{D}, \vec{H}$ through $\vec{E}, \vec{B}$.\\
	From electrodynamics, we know that $\vec{D} := \varepsilon_0\vec{E} + \vec{P}$ \textcolor{gruen}{(in general)} and $\vec{H} := \frac{1}{\mu_0}\vec{B}$ \textcolor{gruen}{(for linear, non-magnetic medii)}, which yields
		\begin{align*}
			\varepsilon_0\vec{\nabla} \vec{E} + \vec{\nabla}\vec{P} &= 0 \\
			\vec{\nabla} \times \vec{E} &= -\partial_t \vec{B} \\
			\vec{\nabla} \vec{B} &= 0 \\
			\frac{1}{\mu_0}\vec{\nabla} \times \vec{B} &= \partial_t (\varepsilon_0 \vec{E} + \vec{P}).
		\end{align*}
	
	Use the nabla theorem $\vec{\nabla} \times (\vec{\nabla} \times \vec{A}) = \vec{\nabla} (\vec{\nabla} \cdot \vec{A}) - \Delta \vec{A}$ on the second equation to decouple $\vec{E}$ and $\vec{B}$:
		\begin{align*}
			\vec{\nabla} \times (\vec{\nabla} \times \vec{E}) &= \vec{\nabla} (\vec{\nabla} \cdot \vec{E}) - \Delta \vec{E} \stackrel{!}{=} -\partial_t (\vec{\nabla} \times \vec{B}) = \\
			&= -\partial_t(\mu_0\partial_t(\varepsilon_0\vec{E} + \vec{P})) = \\
			&= -\mu_0\varepsilon_0 \partial_t^2 \vec{E} - \mu_0 \partial_t^2 \vec{P} \\
			\Leftrightarrow \Delta \vec{E} - \vec{\nabla} (\vec{\nabla} \vec{E}) - \mu_0\varepsilon_0 \partial_t^2 \vec{E} &= \mu_0 \partial_t^2 \vec{P}
		\end{align*}
	
	To simplify this equation, use the approximation, that the electric field \textbf{slowly varies in space and time}. Precisely, that the envelope varies  This is equivalent to $\partial_z |\vec{E}_0| \ll k_z |\vec{E}_0|$ and $\partial_t |\vec{E}_0| \ll \omega |\vec{E}_0|$ \textcolor{gruen}{(the relative change of the envelope is smaller than the relative change of the exponential)}, which lets us neglect the $\vec{\nabla}\vec{E}$ term:
		\begin{align*}
			\Delta \vec{E} - \mu_0\varepsilon_0 \partial_t^2 \vec{E} &= \mu_0 \partial_t^2 \vec{P}
		\end{align*}

	Separate the polarisation into the polarisation of the medium $\vec{P}_M$ and the polarisation of the doped atoms $\vec{P}_A$:
		\begin{align*}
			\Delta \vec{E} - \mu_0\varepsilon_0 \partial_t^2 \vec{E} &= \mu_0 \partial_t^2 (\vec{P}_A + \vec{P}_M) \\
			\Leftrightarrow \Delta \vec{E} - \mu_0\varepsilon_0 \partial_t^2 \vec{E} &= \mu_0 \partial_t^2 (\vec{P}_A + \chi_M\varepsilon_0\vec{E}) \\
			\Leftrightarrow \Delta \vec{E} -  \mu_0\varepsilon_0 \lr{ 1 + \chi_M }\partial_t^2 \vec{E} &= \mu_0 \partial_t^2 \vec{P}_A \\
		\end{align*}

	The factor $1 + \chi_M = \varepsilon_M$ is equal to the refractive index squared $n^2$.\\
	Use the approximation that $|\chi_{M,I}| \ll \chi_{M,R}$. The Taylor expansion of $n$ thus yields:
		\begin{align*}
			n &= \sqrt{ 1 + \chi_M } = \sqrt{ 1 + \chi_{M,R} + i\chi_{M,I}} \approx \sqrt{ 1 + \chi_{M,R} } - \frac{i \chi_{M,I}}{2\sqrt{ 1 + \chi_{M,R}}} \equiv n_R + i\frac{\chi_{M,I}}{2n_R} \\
			\Rightarrow n^2 &= n_R^2 + i\chi_{M,I} + \mathcal{O}(\chi_{M,I}^2)
		\end{align*}
	

\newpage


	Inserting $n^2$ into the wave equation yields:
		\begin{align*}
			\Delta \vec{E} -  \mu_0\varepsilon_0 \lr{ n_R^2 + i\chi_{M,I} }\partial_t^2 \vec{E} &= \mu_0 \partial_t^2 \vec{P}_A
		\end{align*}
	
	Further approximations as well as the use of the ansatz of a monochromatic wave results in the inhomogeneous paraxial wave equation. $\qquad \blacksquare$
\hrule
}



\newpage


\subsection{Semi-classical Description}
The semi-classical description describes
	\begin{enumerate}
		\item light as a classical electromagnetic wave

		\item matter as a system of two-state atoms
	\end{enumerate}


\begin{table}[h]
	\centering
	\caption{Overview results semi-classical description}
	\begin{tabularx}{\textwidth}{X|X}
		\hline\hline
			Successes & Failures \\
		\hline
			- describes the amplification of the E-field & - needs phenomenological damping parameters \\
		\hline
			- describes the effect of a saturation intensity &  \\
		\hline\hline	
	\end{tabularx}
\end{table}


\subsubsection{Maxwell-Bloch Equations}
The equations of motions of the system, describing the
	\begin{itemize}
		\item \textbf{polarization of atomic states}
			\begin{equation}
				\frac{\partial}{\partial t}\vec{P}_0 = \lr{ i(\omega_0 - \omega) - \frac{1}{\tau_P} } \vec{P}_0 + \frac{i\cdot|\vec{d}_{12}|^2}{\hbar} \Delta n \cdot \vec{E}_0,
			\end{equation}

		where $\vec{P}_0 := \varepsilon_0\chi_A \vec{E}_0$.


		\item \textbf{population of the atomic states}
			\begin{equation}
				\frac{\partial}{\partial t} \Delta n = \frac{1}{\hbar} \Im\lrschweif{ \vec{E}_0^* \vec{P}_0 } - \frac{\Delta n - n_0}{\tau}
			\end{equation}


		\item \textbf{electric field}
			\begin{equation}
				\lr{ \frac{\partial}{\partial z} + \frac{1}{c_g}\frac{\partial}{\partial t} }\vec{E}_0 = \frac{i}{2}\mu_0 \omega c_p \vec{P}_0
			\end{equation}
	\end{itemize}

are called \textbf{Maxwell-Bloch equations} and include the \textbf{population inversion} $\Delta n := n_2 - n_1$ \textcolor{gruen}{(and the population densities $n_i = |c_i|^2$)}, \textbf{transition frequency} (between the two states) $\omega_0$, \textbf{damping time} $\tau_P$ \textcolor{gruen}{(phenomenological)}, group velocity $c_g$, phase velocity $c_p$.


\newpage


{\color{violet}
\hrule
	\fbox{Derivation}\\
	Strategy:\\
	1) Determine the occupation coefficients $c_1$ and $c_2$ for the two-state atom\\
	2) Using $c_1$ and $c_2$, calculate the polarisation \\
	3) Calculate the dynamic of the polarisation and the population inversion \\
	4) Add phenomenological damping parameters \\
	\\
	zz: Equation of motion for the population of the atomic states \\
	\\
	Consider the Hamiltonian of the semi-classical system under \textbf{slowly varying E-field approximation}
		\begin{align*}
			\hat{H} &= \hat{H}_0 + \hat{H}_\text{int} = \\
			&= \hat{H}_0 + \hat{\vec{d}}\vec{E}(\vec{x}, t) \approx \\
			&\approx \hat{H}_0 + \hat{\vec{d}}\vec{E}(t).
		\end{align*}
	
	with the dipole moment operator $\hat{\vec{d}} = -e\vec{x}$.\\
	\\
	As an ansatz for the wave function, consider
		\begin{align*}
			\ket{\psi(t)} &= C_1(t) \ket{\phi_1} + C_2(t) \ket{\phi_2} = \\
			&\equiv c_1(t) e^{-i\omega_1 t} \ket{\phi_1} + c_2(t) e^{-i\omega_2 t} \ket{\phi_2}.
		\end{align*}
	
	To determine the time evolution, plug both into the time dependent Schrödinger equation:
		\begin{align*}
			i\hbar \partial_t \ket{\psi(t)} &= \hat{H} \ket{\psi(t)} \\
			\Leftrightarrow i\hbar \partial_t \lr{ c_1(t) e^{-i\omega_1 t} \ket{\phi_1} + c_2(t) e^{-i\omega_2 t} \ket{\phi_2} } &= \lr{ \hat{H}_0 + \hat{H}_\text{int} } [ c_1(t) e^{-i\omega_1 t} \ket{\phi_1} + \\
			&\phantom{=}+ c_2(t) e^{-i\omega_2 t} \ket{\phi_2} ] \\
			\Leftrightarrow i\hbar \lr{ \dot{c_1} e^{-i\omega_1 t} - i\omega_1 c_1 e^{-i\omega_1 t} } \ket{\phi_1} + i\hbar \lr{ \dot{c_2} e^{-i\omega_2 t} - i\omega_2 c_2 e^{-i\omega_2 t} } \ket{\phi_2} &= \lr{ c_1 e^{-i\omega_1 t} E_1 + \hat{H}_\text{int} c_1 e^{-i\omega_1 t} } \ket{\phi_1} + \\
			&\phantom{=} + \lr{ c_2 e^{-i\omega_2 t} E_2 + \hat{H}_\text{int} c_2 e^{-i\omega_2 t} } \ket{\phi_2},
		\end{align*}

	where in the last step, the eigenvalue equations $\hat{H}_0 \ket{\phi_i} = E_i \ket{\phi_i}$ with $i\in\{1,2\}$ were used.\\
	\\
	Multiply the last equation with $\bra{\phi_1}$ and $\bra{\phi_2}$ respectively.\\
	Begin with $\bra{\phi_1}$:
		\begin{align*}
			i\hbar \lr{ \dot{c_1} e^{-i\omega_1 t} - i\omega_1 c_1 e^{-i\omega_1 t} } + 0 &= \lr{ c_1 e^{-i\omega_1 t} E_1 + c_1 e^{-i\omega_1 t} \underbrace{ \braket{ \phi_1 | \hat{H}_\text{int} | \phi_1 }}_{=0} } + \lr{ 0 + c_2 e^{-i\omega_2 t} \braket{ \phi_1 | \hat{H}_\text{int} | \phi_2 } } \\
			\Leftrightarrow i\hbar \dot{c}_1 e^{-i\omega_1 t} + \cancel{ \underbrace{ \hbar\omega_1 }_{\comment{= E_1}} c_1 e^{-i\omega_1 t} } &= \cancel{ c_1 e^{-i\omega_1 t} E_1 } + c_2 e^{-i\omega_2 t} \braket{ \phi_1 | \hat{H}_\text{int} | \phi_2 } \\
			\stackrel{\comment{\cdot e^{+i\omega_1 t}}}{\Leftrightarrow} i\hbar \dot{c_1} &= c_2 e^{-i(\omega_2 - \omega_1) t} \braket{ \phi_1 | \hat{H}_\text{int} | \phi_2 }
		\end{align*}

	Define the \textbf{dipole matrix element} $\vec{d}_{ij} := \braket{ \phi_i | (-e\vec{x}) | \phi_j }$ as well as the \textbf{transition frequency} $\omega_0 := \omega_2 - \omega_1$, which simplifies the equation above to
		\begin{align*}
			i\hbar \dot{c}_1 &= c_2 e^{-i\omega_0 t} \vec{d}_{12}\vec{E}(t).
		\end{align*}
	
	Multiplication with $\bra{\phi_2}$ yields analogously
		\begin{align*}
			i\hbar\dot{c}_2 &= c_1 e^{+i\omega_0 t} \vec{d}_{12}^\star \vec{E}(t).
		\end{align*}


	An expression, that will turn up later, is the time derivative of the product of $c_1$ and $c_2$:
		\begin{align*}
			\partial_t (c_1^\star c_2) &= \dot{c}_1^\star c_2 + c_1^\star \dot{c}_2 = c_2 \lr{ -\frac{1}{i\hbar} c_2 e^{+i\omega_0 t} \vec{d}_{12} \vec{E}(t) } + c_1^\star \lr{ \frac{1}{i\hbar} e^{i\omega_0 t} \vec{d}_{12} \vec{E}(t) } = \\
			&= \frac{1}{i\hbar} \vec{d}_{12}^\star \vec{E} \lr{ |c_1|^2 - |c_2|^2 } e^{i\omega_0 t}
		\end{align*}
	
	Now, introduce the \textbf{macroscopic polarisation}
		\begin{align*}
			\vec{P}(t) := n_0 \cdot \braket{ \psi(t) | \vec{d} | \psi(t) }
		\end{align*}
	
	which becomes in this two-state atom system
		\begin{align*}
			\vec{P} &= n_0 \lr{ c_1^\star e^{+i\omega_1 t} \bra{\phi_1} + c_2^\star e^{+i\omega_2 t} \bra{\phi_2} } \vec{d} \lr{ c_1 e^{-i\omega_1 t} \ket{\phi_1} + c_1^\star e^{-i\omega_2 t} \ket{\phi_2} } = \\
			&= n_0 \lr{ 0 + c_1^\star c_2 e^{-i(\omega_2 - \omega_1)t} \braket{ \phi_1 | \vec{d} | \phi_2 } + c_1 c_2^\star e^{-i(\omega_1 - \omega_2)t} \braket{ \phi_2 | \vec{d} | \phi_1} + 0 } \equiv \\
			&\equiv \vec{P}^+ + \vec{P}^-,
		\end{align*}
	
	with the definitions
		\begin{align*}
			\vec{P}^+ &= n_0 c_1^\star c_2 e^{-i\omega_0 t} \vec{d}_{12} \\
			\vec{P}^- &= n_0 c_1 c_2^\star e^{+i\omega_0 t} \vec{d}_{12}^\star.
		\end{align*}
	
	Now, after these auxiliary calculations, we can consider the equation of motion for the \textbf{population inversion} $\Delta n := n_2 - n_1$.\\
	Calculate $\partial_t n_2$ and $\partial_t n_1$ seperately:
		\begin{align*}
			\bullet \ \partial_t n_2 &= \partial_t (n_0 |c_2|^2) = n_0 \dot{c}_2^\star c_2 + n_0 c_2^\star \dot{c}_2 = \\
			&= n_0 c_2 \lr{ -\frac{1}{i\hbar} \vec{d}_{12} \vec{E} e^{+i\omega_0 t} } + n_0 c_2^\star \lr{ \frac{1}{i\hbar} \vec{d}_{12}^\star \vec{E} c_1 e^{-i\omega_0 t} } c_2^\star \equiv \\
			&\equiv -\frac{1}{i\hbar} \vec{P}^+ \vec{E} + \frac{1}{i\hbar} \vec{E}\vec{P}^- \approx \\
			&\stackrel{\commenttext{RWA}}{\approx} -\frac{1}{i\hbar} \vec{P}^+ \vec{E}^- + \frac{1}{i\hbar} \vec{E}^+\vec{P}^- = -\frac{1}{i\hbar} \lr{ \vec{E}^- \vec{P}^+ - \vec{E}^+ \vec{P}^- } = \\
			&\stackrel{\textcolor{red}{Wieso +?}}{=} \frac{2}{\hbar} \Im \lrschweif{ \vec{E}^-\vec{P}^+ } \equiv \\
			&\equiv \frac{2}{\hbar} \Im \lrschweif{ \lr{ E_0^\star e^{+i\omega t} } \lr{ \frac{1}{2} \vec{P}_0 e^{-i\omega t} } } = \frac{1}{2} \frac{1}{\hbar} \Im \lrschweif{ \vec{E}_0^\star \vec{P}_0 } \\
			\bullet \ \partial_t n_1 &= ... = - \frac{1}{2} \frac{1}{\hbar} \Im \lrschweif{ \vec{E}_0^\star \vec{P}_0 } \\
			\Rightarrow \partial_t \Delta n &= \partial_t n_2 - \partial_t n_1 = \frac{1}{\hbar}  \Im \lrschweif{ \vec{E}_0^\star \vec{P}_0 } \lr{ \frac{1}{2} - \lr{-\frac{1}{2}} } = \frac{1}{\hbar} \Im \lrschweif{ \vec{E}_0^\star \vec{P}_0 }
		\end{align*}

	Finally, add the phenomenological term $-\frac{\Delta n - n_0}{\tau}$ to arrive at the desired equation
		\begin{align*}
			\partial_t \Delta n &= \frac{1}{\hbar} \Im \lrschweif{ \vec{E}_0^\star \vec{P}_0 } - \frac{\Delta n - n_0}{\tau}. \qquad \blacksquare
		\end{align*}

	
	\fbox{Derivation}\\
	zz: Equation of motion for the polarization fo the atomic states\\
	\\
	Calculate the partial derivative of $\vec{P}^\pm$ respective the time $t$:
		\begin{align*}
			\bullet \ \partial_t \vec{P}^+ &= \partial_t \lr{ n_0 c_1^\star c_2 e^{-i\omega_0 t} \vec{d}_{12} } = n_0 \partial_t\lr{ c_1^\star c_2 } e^{-i\omega_0 t} \vec{d}_{12} + (-i\omega_0) n_0 c_1^\star c_2 e^{-i\omega_0 t} \vec{d}_{12}
		\intertext{Insert $\partial_t(c_1^\star c_2)$, as calculated in the derivation of the equation of motion for the atomic population, and recognize $\vec{P}^+$ in the second term:}
			&= n_0 \lreck{ \frac{\vec{d}_{12}^\star \vec{E}}{i\hbar} \lr{ |c_1|^2 - |c_2|^2 } e^{+i\omega_0 t} } e^{-i\omega_0 t} \vec{d}_{12} -i\omega_0 \vec{P}^+ =
		\intertext{Use, that $\vec{d}_{12}$ and $\vec{E}$ point in the same direction \textcolor{red}{(Warum?)}, which yields $(\vec{d}_{12}^\star \vec{E}) \vec{d}_{12} = |\vec{d}_{12}|^2 \vec{E}$ as well as $n_i = n_0 |c_i|^2$:}
			&= \frac{|\vec{d}_{12}|^2 \vec{E}}{i\hbar} ( n_1 - n_2) - i\omega_0 \vec{P}^+ =
		\intertext{Finally, insert the \textbf{phenomenlogical damping term} $-\vec{P}^+ / \tau_P$ and recognise the population inversion $\Delta n$}
			&= \lr{ -i\omega_0 - \frac{1}{\tau_P} } \vec{P}^+ + i\frac{|\vec{d}_{12}|^2 \vec{E}}{\hbar} \Delta n.
		\end{align*}
	
	Completely analogous, it follows that
		\begin{align*}
			\bullet \ \partial_t \vec{P}^- &= \lr{ + i\omega_0 - \frac{1}{\tau_P} } \vec{P}^- + i\frac{|\vec{d}_{12}|^2\vec{E}}{\hbar} \Delta n.
		\end{align*}
	
	By inserting $\vec{P}^\pm = \vec{P}_0 e^{\mp\omega t}$ \textcolor{red}{(Woher kommt das?)} into both equations, one finds one equation of motion for $\vec{P}_0$:
		\begin{align*}
			\partial_t \vec{P}_0 &= \lr{ i(\omega - \omega_0) } \vec{P}_0 + i\frac{|\vec{d}_{12}|^2 \vec{E}}{\hbar} \Delta n \qquad \blacksquare
		\end{align*}

	\fbox{Derivation}\\
	zz: Equation of motion for the electric field\\
	See classical description of light-matter interaction for the case of no losses $\alpha = \kappa = 0$. $\qquad \blacksquare$\\
	\textcolor{red}{Wieso  $\alpha = \kappa = 0$?}
\hrule
}


\newpage


\subsubsection*{Coherent Regime}
In the coherent regime, the damping terms $\frac{1}{\tau_P}$ and $-\frac{\Delta n - n_0}{\tau}$ in the Maxwell-Bloch equations are neglected.\\
The Maxwell-Bloch equations can be written as
	\begin{equation}
		\frac{\diff}{\diff t}\vec{\mu} = \vec{\mu} \times \vec{R} = 
		\begin{pmatrix}
			\frac{\Re P_0}{n_0 \cdot |d_{12}|} \\
			\frac{\Im P_0}{n_0 \cdot |d_{12}|} \\
			\frac{\Delta n}{n_0}
		\end{pmatrix}
		\times
		\begin{pmatrix}
			\Re \Omega_R \\
			\Im \Omega_R \\
			-\Delta
		\end{pmatrix}
	\end{equation}

with the \textbf{Rabi-frequency} $\Omega_R := \frac{|d_{12}|}{\hbar}E_0$ and the \textbf{detuning} $\Delta := \omega - \omega_0$.


\subsubsection*{Incoherent Regime}
In the incoherent regime, the time derivates $\partial_t \vec{P}_0$ is neglected:
	\begin{equation}
		0 = \lr{ i(\omega_0 - \omega) - \frac{1}{\tau_P} } \vec{P}_0 + \frac{i\cdot|\vec{d}_{12}|^2}{\hbar} \Delta n \cdot \vec{E}_0
	\end{equation}

Thus, the polarization $\vec{P}$ instantaneously follows the electric field $\vec{E}$.


\subsubsection{Absorption, Amplification and Saturation Intensity}
The comparison between the semi-classically and classically derived susceptibility yields the semi-classical gain coefficient, which depends on the population inversion $\Delta n$:
	\begin{equation}
		g(\Delta n,\omega) = \Delta n \cdot \sigma(\omega)
		\Bigg\{
			\begin{matrix}
				<0, & \Delta n<0, & \text{absorption} \\
				>0, & \Delta n>0, & \text{amplification} 
			\end{matrix}
	\end{equation}

with the cross section $\sigma(\omega)$.\\
\\
A consequence of the 3rd MBE is the intensity dependend and restricted population inversion and thus gain factor
	\begin{equation}
		\Delta n = \frac{n_0}{ 1 + \frac{2\sigma(\omega)\tau}{\hbar\omega} I } \equiv \frac{n_0}{ 1 + \frac{I}{I_\text{sat}} }.
	\end{equation}


\newpage


{\color{violet}
\hrule
	\fbox{Herleitung}\\
	zz: Absorption coefficient $\alpha$\\
	\\
	The first Maxwell-Bloch equation in the incoherent regime reads
		\begin{align*}
			0 &= \lr{ i(\omega_0 - \omega) - \frac{1}{\tau_P} } \vec{P}_0 + \frac{i\cdot|\vec{d}_{12}|^2}{\hbar} \Delta n \cdot \vec{E}_0 \\
			\Leftrightarrow \vec{P}_0 &= -\frac{i \cdot |\vec{d}_{12}|^2}{\hbar} \Delta n \cdot \frac{1}{ \lr{ i(\omega_0 - \omega) - \frac{1}{\tau_P} } }\cdot \vec{E}_0.
		\end{align*}

	Since a linear medium is assumed in the semi-classical description of light-matter interaction, the relationship between polarisation and electric field is
		\begin{align*}
			\vec{P}_0 &= \varepsilon_0 \chi_A \cdot \vec{E}_0.
		\end{align*}
	
	Comparing both equations yields the semi-classical susceptibility
		\begin{align*}
			\chi_A &= \Delta n \cdot \frac{ |\vec{d}_{12}|^2 / (\varepsilon_0 \hbar) }{(\omega_0 - \omega) + \frac{i}{\tau_P}} 
		\end{align*}

	Finally, we compare this susceptibility with the classical one, which introduces the absorption coefficient $\alpha$:
		\begin{align*}
			\Im \lrschweif{ \Delta n \cdot \frac{ |\vec{d}_{12}|^2 / (\varepsilon_0 \hbar) }{(\omega_0 - \omega) + \frac{i}{\tau_P}} } &\stackrel{!}{=} \frac{\alpha}{k_0} \\
			\Leftrightarrow \Im \lrschweif{ \Delta n \cdot \frac{ |\vec{d}_{12}|^2 / (\varepsilon_0 \hbar) }{ (\omega_0 - \omega)^2 + \frac{1}{\tau_P^2} } \lr{ (\omega_0 - \omega) - \frac{i}{\tau_P} } } &= \frac{\alpha}{k_0} \\
			\Leftrightarrow -\Delta n \cdot \frac{ |\vec{d}_{12}|^2 / (\varepsilon_0 \hbar) }{(\omega_0 - \omega)^2 + \frac{1}{\tau_P^2}} \cdot \frac{1}{\tau_P} &= \frac{\alpha}{k_0} \\
			\Leftrightarrow \alpha(\Delta n) &= -\Delta n \cdot \frac{ |\vec{d}_{12}|^2 / (\varepsilon_0 \hbar) }{(\omega_0 - \omega)^2 + \frac{1}{\tau_P^2}} \cdot \frac{1}{\tau_P} \equiv \\
			&\equiv -\Delta n \cdot \sigma(\omega) \qquad \blacksquare
		\end{align*}


	\fbox{Derivation}
	zz: $\Delta n = -\frac{n_0}{ 1 + \frac{I}{I_\text{sat}} }$\\
	\\
	Consider the 2nd MBE in the incoherent regime
		\begin{align*}
			0 &= \frac{1}{\hbar} \Im\{ \vec{E}_0^\star \vec{P}_0 \} - \frac{\Delta n - n_0}{\tau}
		\intertext{Insert $\vec{P_0} = \varepsilon_0 \chi_{A,I} \vec{E}_0$:}
			0 &= \frac{1}{\hbar} |\vec{E}_0|^2 \varepsilon_0 \chi_{A,I} - \frac{\Delta n - n_0}{\tau} \\
			\intertext{Insert the intensity $I = \frac{1}{2} \varepsilon_0 |\vec{E}_0|^2 c$}
			0 &= \frac{1}{\hbar} \frac{2}{c} I \chi_{A,i} - \frac{\Delta n - n_0}{\tau}
			\intertext{Rearrange $\alpha = k \chi_{A,i} = -\Delta n \sigma(\omega)$ for $\chi_{A,i}$ and insert it into the equation:}
			0 &= - \frac{1}{\hbar} \frac{2}{c} I \lr{ -\frac{\Delta n \sigma(\omega)}{k} } - \frac{\Delta n - n_0}{\tau} \\
			0 &= + \frac{1}{\hbar} \frac{2}{\omega} I \Delta n \sigma(\omega) - \frac{\Delta n - n_0}{\tau} \\
			\Leftrightarrow 0 &= -\Delta n \lr{ \frac{1}{\hbar} \frac{2}{\omega} I \sigma(\omega) + \frac{1}{\tau} } + \frac{n_0}{\tau} \\
			\Leftrightarrow \Delta n &= \frac{n_0}{ \frac{2\tau \sigma(\omega)}{\hbar\omega} I + 1 } \equiv \frac{n_0}{1 + \frac{I}{I_\text{sat}}} \qquad\blacksquare
		\end{align*}

\hrule
}


\newpage


\subsubsection{Spontaneous Emission, Stimulated Emission and Absorption}
\begin{figure}[h]
	\centering
	\fbox{
	\includegraphics[width=0.6\textwidth]{Grafiken/spontaneous emission.pdf}
	}
	\caption{Figure of spontaneous emission, stimulated emission and absorption with their respective Einstein coefficients}
\end{figure}

The transition rates of these processes are describes by the \textbf{Einstein coefficients}
	\begin{align*}
		A &= \frac{2}{3} \frac{\omega^3}{h \varepsilon_0 c^3} \cdot |\vec{d}_{12}|^2 \\
		B_{21} \cdot u(\omega_{21}) &= \frac{2}{3} \frac{\pi^2}{h^2 \varepsilon_0} \cdot |d_{12}|^2 \cdot u(\omega_{21}) \\
		B_{12} \cdot u(\omega_{21}) &= B_{21} \cdot u(\omega_{21}).
	\end{align*}

The relation between the coefficients is
	\begin{align*}
		A &= \frac{\hbar \omega_{21}^3}{\pi^2 c^3} \cdot B.
	\end{align*}

	\textbf{Phenomenologically, one finds a connection between the rates described by the Einstein coefficients and the cross section $\sigma(\nu)$.}


\newpage


{\color{violet}
\hrule
	\fbox{Derivation}\\
	ges: Einstein coefficients $B_{21}, B_{12}$\\
	\\
	The Einstein coefficients can be derived from \underline{Fermi's Golden rule}
		\begin{align*}
			P_{n\to k} &= |a_k^{(1)}|^2 = \frac{ | \braket{ k | \hat{H}' | n } |^2 }{\hbar^2} \cdot \frac{ \sin^2\lr{ \frac{\omega_{kn}}{2} T } }{ \lr{ \frac{\omega_{kn}}{2} }^2 },
		\end{align*}
	
	where $\omega_{kn} := \omega_k - \omega_n$.\\
	\\
	In the case of a periodic perturbation 
		\begin{align*}
			\hat{H}' &= \vec{d} \cdot \vec{E} = \vec{d} \cdot \vec{E}_0 \cdot \cos(\omega t) = \vec{d} \cdot \frac{\vec{E}_0}{2} \lr{ e^{i\omega t} + e^{-i\omega t} } = \\
			&= d \cdot \frac{E_0}{2} \lr{ e^{i\omega t} + e^{-i\omega t} }
		\end{align*}

	with the dipole moment $\vec{d}$, the formula becomes (see advanced quantum mechanics lecture):
		\begin{align*}
			P_{n\to k} &= \frac{2\pi}{\hbar} \cdot T \cdot \lreck{ |\braket{ k | \frac{d \cdot E_0}{2} | n }|^2 \cdot \delta ( E_k - E_n - \hbar\omega ) + | \braket{ k | \frac{d \cdot E_0}{2} | n }|^2 \cdot \delta ( E_k - E_n + \hbar\omega ) } = \\
			&= \frac{2\pi}{\hbar} \cdot T \cdot \frac{E_0^2}{4} \cdot \lreck{ |\braket{ k | d | n }|^2 \cdot \delta ( E_k - E_n - \hbar\omega ) + | \braket{ k | d | n }|^2 \cdot \delta ( E_k - E_n + \hbar\omega ) } \equiv \\
			&\equiv \frac{\pi \cdot E_0^2 \cdot |d_{kn}|^2}{2\hbar} \cdot T \cdot \lreck{ \delta ( E_k - E_n - \hbar\omega ) + \delta ( E_k - E_n + \hbar\omega ) } = \\
			&= \frac{\pi \cdot E_0^2 \cdot |d_{kn}|^2}{2\hbar^2} \cdot T \cdot \lreck{ \delta ( \omega_k - \omega_n - \omega ) + \delta ( \omega_k - \omega_n + \omega ) }
		\end{align*}

	For a \underline{two-state atom}, the transition probability become:
		\begin{align*}
			P_{1\to 2} &= \frac{\pi \cdot E_0^2 \cdot |d_{21}|^2}{2\hbar^2} \cdot T \cdot \delta(\omega_2 - \omega_1 - \omega) \\
			P_{2\to 1} &= \frac{\pi \cdot E_0^2 \cdot |d_{12}|^2}{2\hbar^2} \cdot T \cdot \delta(\omega_1 - \omega_2 + \omega)
		\end{align*}
	
	Thus, the transition rates are:
		\begin{align*}
			\Gamma_{1\to 2} &= \frac{\pi \cdot E_0^2 \cdot |d_{21}|^2}{2\hbar^2} \cdot \delta(\omega_2 - \omega_1 - \omega) \\
			\Gamma_{2\to 1} &= \frac{\pi \cdot E_0^2 \cdot |d_{12}|^2}{2\hbar^2} \cdot \delta(\omega_1 - \omega_2 + \omega)
		\end{align*}

	For \underline{non-monochromatic light}, one ''adds'' the rates of every frequency. There, the amplitudes can vary from frequency to frequency. The distribution of amplitudes is described by $u(\omega) = \frac{\varepsilon_0 \cdot E_0(\omega)^2}{2}$:
		\begin{align*}
			\Gamma_{1\to 2}^\text{total} &= \int_0^\infty \diff\omega \, \frac{\pi \cdot |d_{21}|^2}{\varepsilon_0\hbar^2} \cdot u(\omega) \cdot \delta(\omega_2 - \omega_1 - \omega) = \\
			&= \frac{\pi \cdot |d_{21}|^2}{\varepsilon_0\hbar^2} \cdot u(\omega_{21}) \equiv\\
			&\equiv B_{12} \cdot \omega_{21} \\
			\Gamma_{2\to 1}^\text{total} &= \int_0^\infty \diff\omega \, \frac{\pi \cdot |d_{21}|^2}{\varepsilon_0\hbar^2} \cdot u(\omega) \cdot \delta(\omega_1 - \omega_2 + \omega) = \\
			&= \frac{\pi \cdot |d_{12}|^2}{\varepsilon_0\hbar^2} \cdot u(\omega_{12}) \equiv \\
			&\equiv B_{12} \cdot u(\omega_{12})
		\end{align*}
	
	\textcolor{blue}{In the case of non-perpendicular $\vec{d}$ and $\vec{E}$ vectors, the transition rate depends on the angle $\theta$ between both vectors. Averageing over these vectors yields the factor $\frac{2}{3}$.}\\
	\\
	\fbox{Derivation}\\
	ges: Einstein coefficient $A$\\
	\\
	The Einstein coefficient for spontaneous emission $A$ has to be derived using quantum electrodynamics.\\
	One can however use Einstein's thermodynamical argument.\\
	\\
	1) Formulate rate equations\\
	2) Calculate $N_2/N_1$ and set it equal to the Boltzmann distribution\\
	3) Use $\rho\to\infty$ for $T\to\infty$ to find $B$ coefficients \\
	4) Consider small frequencies \\
	5) Compare with Rayleigh-Jeans law\\
	\\
	\textcolor{red}{ERGÄNZEN}
\hrule
}


\newpage


\subsection{Quantum Mechanical Description}
The quantum mechanical description describes
	\begin{enumerate}
		\item light as a quantized electromagnetic wave
		\item matter as two-state atoms\\
		A typical assumptions (\textbf{laser medium}): linear, isotropic, non-magnetic
	\end{enumerate}


\begin{table}[h]
	\centering
	\caption{Overview results quantum mechanical description}
	\begin{tabularx}{\textwidth}{X|X}
		\hline\hline
			Successes & Failures \\
		\hline
			 &  \\
		\hline
			&  \\
		\hline\hline	
	\end{tabularx}
\end{table}


\newpage


\subsubsection{Quantization of the Electromagnetic Field}
\subsubsection*{Single-Mode Field}
The quantization of the electromagnetic field yields the following expressions for the E- and B-field and the Hamiltonian
	\begin{align}
		\hat{\vec{E}}_n(z,t) &= \vec{E}_0 \cdot \lr{ \hat{a}_n(0) \cdot e^{-i\omega_n t} + \hat{a}_n^\dagger(0) \cdot e^{+i\omega t} } \cdot \sin(k_n z) \\
		\hat{\vec{B}}_n(z,t) &= \frac{\vec{B}_0}{i} \cdot \lr{ \hat{a}_n(0) e^{-i\omega_n t} - \hat{a}_n^\dagger(0) \cdot e^{+i\omega_n t} } \cdot \sin(k_n z) \\
		\vec{H}_n &= \hbar\omega \lr{ \hat{a}_n^\dagger \hat{a}_n + \frac{1}{2} },
	\end{align}

where $k_n := \frac{\pi}{L} n$ with $n = 1,2,...$ and the ladder operators
	\begin{align}
		\hat{a}_j \ket{\tilde{n}} &= \sqrt{\tilde{n}} \cdot \ket{\tilde{n}-1} \\
		\hat{a}_j^\dagger \ket{\tilde{n}} &= \sqrt{\tilde{n}+1} \cdot \ket{\tilde{n}+1} \\
		\hat{a}_j \ket{0} &= 0,
	\end{align}

where the kets are eigenstates of the Hamiltonian.


\subsubsection*{Multi-Mode Field}
For multi-mode fields, the second quantization yields:
	\begin{align}
		\hat{\vec{E}}(\vec{x},t) &= i \sum_{\vec{k},s} \sqrt{ \frac{\hbar \omega_k }{ 2\varepsilon_0 V } } \vec{e}_{\vec{k},s} \lr{ \hat{a}_{\vec{k},s}(0) \cdot e^{i(\vec{k}\vec{r} - \omega_k t)} - \hat{a}_{\vec{k},s}^\dagger(0) \cdot e^{-i(\vec{k}\vec{r} - \omega_k t)} } \\
		\hat{\vec{B}}(\vec{x},t) &= \frac{i}{c} \sum_{\vec{k},s} \sqrt{ \frac{\hbar \omega_k }{ 2\varepsilon_0 V } } \lr{ \frac{\vec{k}}{|\vec{k}|} \times \vec{e}_{\vec{k},s} } \cdot \lr{ \hat{a}_{\vec{k},s}(0) \cdot e^{i(\vec{k}\vec{r} - \omega_k t)} - \hat{a}_{\vec{k},s}^\dagger(0) \cdot e^{-i(\vec{k}\vec{r} - \omega_k t)} } \\
		\vec{H} &= \sum_{\vec{k},s} \hbar\omega_{\vec{k}}\lr{ \hat{a}_{\vec{k},s}^\dagger \hat{a}_{\vec{k},s} + \frac{1}{2} }.
	\end{align}

with the ladder operators
	\begin{align}
		\hat{a}_j \ket{ n_1, n_2, ..., n_j, ...} &= \sqrt{n_j} \cdot \ket{ n_1, n_2, ..., n_j - 1, ...} \\
		\hat{a}_j^\dagger \ket{ n_1, n_2, ..., n_j, ...} &= \sqrt{n_j + 1} \cdot \ket{ n_1, n_2, ..., n_j + 1, ...} \\
		\hat{a}_j \ket{..., 0_j, ...} &= 0,
	\end{align}

where the kets are eigenstates of the Hamiltonian.


\newpage


{\color{violet}
\hrule
	\fbox{Derivation Single-Mode Quantization}\\
	Strategy:\\
	1) Formulate the E- and B-field classically \underline{with} boundary condition \\
	2) Formulate the classical Hamilton function\\
	3) Recognize terms in the Hamiltonian as position and momentum\\
	4) Replace position and momentum with their respective operators $\hat{p}, \hat{q}$\\
	5) Express $\hat{p}, \hat{q}$ with the ladder operators $\hat{a}, \hat{a}^\dagger$\\
	\\
	Consider an electromagnetic field between two mirrors, which are located at $z=0$ and $z=L$.\\
	The Maxwell equations in the vacuum are
		\begin{align*}
			\vec{\nabla} \cdot \vec{E} &= 0 \\
			\vec{\nabla} \cdot \vec{B} &= 0 \\
			\vec{\nabla} \times \vec{E} &= - \partial_t \vec{B} \\
			\vec{\nabla} \times \vec{B} &= \varepsilon_0\mu_0 \partial_t \vec{E}.
		\end{align*}
	
	The boundary conditions enforce
		\begin{align*}
			\vec{E}(z=0,t) &= \vec{E}(z=L,t) = \vec{0}.
		\end{align*}
	
	Two solution to the Maxwell equations together with the boundary conditions are
		\begin{align*}
			\vec{E}(z,t) &= \sqrt{ \frac{2\omega_n^2}{V\varepsilon_0} } \cdot q_n(t) \cdot \sin(k_n z) \cdot \vec{e}_x \\
			\vec{B}(z,t) &= \frac{ \mu_0 \varepsilon_0 }{k_n} \sqrt{ \frac{2\omega_n^2}{V\varepsilon_0} } \cdot \dot{q}_n(t) \cdot \cos(k_n z) \cdot \vec{e}_y,
		\end{align*}

	with the amplitude $q_n(t)$ and where $k_n := \frac{\pi}{L} \cdot n$, $\omega_n = c \cdot k_n$ and the mode index $n = 1,2,...$.\\
	\\
	The Hamilton funcion is the total energy of the system, so in this case the integral of the electric and magnetic energy densities from $z=0$ to $z=L$:
		\begin{align*}
			H &= \int_0^L \diff z \, \lr{ \frac{\varepsilon_0 }{2} \vec{E}^2(z,t) + \frac{1}{2\mu_0} \vec{B}(z,t)^2 } = \\
			&= \int_0^L \diff z \, \lreck{ \frac{\varepsilon_0}{2} \lr{ \sqrt{ \frac{2\omega_n^2}{V\varepsilon_0} } \cdot q_n(t) \cdot \sin(k_n z) \cdot \vec{e}_x }^2 + \frac{1}{2\mu_0} \lr{ \frac{ \mu_0 \varepsilon_0 }{k_n} \sqrt{ \frac{2\omega_n^2}{V\varepsilon_0} } \cdot \dot{q}_n(t) \cdot \cos(k_n z) \cdot \vec{e}_y }^2 } = \\
			&= \int_0^L \diff z \, \lreck{ \frac{\varepsilon_0}{2} \cdot \frac{2\omega_n^2}{V\varepsilon_0} \cdot q_n^2(t) \cdot \sin_n^2(k_n z) + \frac{1}{2\mu_0} \cdot \frac{\mu_0^2\varepsilon_0^2}{k_n^2} \cdot \frac{2\omega_n^2}{V\varepsilon_0} \cdot \dot{q}_n^2(t) \cdot \cos^2(k_n z) } = \\
			&= \int_0^L \diff z \, \lr{ \frac{\omega_n^2}{V} \cdot q_n^2(t) \cdot \sin_n^2(k_n z) + \underbrace{ \frac{\mu_0\varepsilon_0}{k_n^2} }_{= 1 / \omega_n^2}\cdot \frac{\omega_n^2}{V} \cdot \dot{q}^2_n(t) \cdot \cos^2(k_n z) } = \\
			&= \frac{\omega_n^2}{V} \cdot q_n^2(t) \cdot \int_0^L \diff z \, \sin_n^2(k_n z) + \frac{1}{V} \cdot \dot{q}_n(t) \cdot \int_0^L \diff z \, \cos^2(k_n z) = \\
			&= \frac{\omega_n^2}{V} \cdot q_n^2(t) \cdot \frac{V}{2} + \frac{1}{V} \cdot \dot{q}^2_n(t) \cdot \frac{V}{2} = \\
			&= \frac{1}{2} \lreck{ \omega_n^2 \cdot q_n^2(t) + \dot{q}^2_n(t) } = \\
			&= \frac{1}{2} \lreck{ \omega_n^2 \cdot q_n^2(t) + p_n^2(t) }
		\end{align*}

	Recognise this Hamilton function as the Hamilton function of the harmonic oscillator.\\
	Thus, $q_n(t)$ and $p_n(t)$ can be interpreted as the position and momentum.\\
	\\
	Replace $q_n(t)$ and $p_n(t)$ with their respective operators $\hat{q}_n(t)$ and $\hat{p}_n(t)$.\\
	\\
	The relation between these operators and the ladder operators are
		\begin{align*}
			\hat{q}_n(t) &= \sqrt{ \frac{\hbar}{2\omega_n} } \lr{ \hat{a}_n + \hat{a}^\dagger_n } \\
			\hat{p}^\dagger_n(t) &= -i\sqrt{ \frac{\hbar\omega_n}{2} } \lr{ \hat{a}_n - \hat{a}^\dagger_n }.
		\end{align*}

	Inserting these back into E- and B-field yields:
		\begin{align*}
			\hat{\vec{E}}(z,t) &= \sqrt{ \frac{2\omega_n^2}{V\varepsilon_0} } \cdot \sqrt{ \frac{\hbar}{2\omega_n} } \lr{ \hat{a}_n + \hat{a}^\dagger_n } \cdot \sin(k_n z) \cdot \vec{e}_x = \\
			&= \sqrt{ \frac{ \hbar\omega_n }{V \varepsilon_0 } } \cdot ( \hat{a}_n + \hat{a}^\dagger ) \cdot \sin(k_n z) \cdot \vec{e}_x \equiv \\
			&\equiv \vec{E}_0 \cdot ( \hat{a}_n + \hat{a}^\dagger ) \cdot \sin(k_n z) \qquad \blacksquare \\
			\hat{\vec{B}}(z,t) &= \frac{ \mu_0 \varepsilon_0 }{k_n} \sqrt{ \frac{2\omega_n^2}{V\varepsilon_0} } \cdot \lr{ -i\sqrt{ \frac{\hbar\omega_n}{2} } \lr{ \hat{a}_n - \hat{a}^\dagger_n } } \cdot \cos(k_n z) \cdot \vec{e}_y \equiv \\
			&\equiv \frac{B_0}{i} \cdot (\hat{a}_n - \hat{a}^\dagger) \cdot \sin(k_n z) \qquad \blacksquare
		\end{align*}
	
	The time evolution of the ladder operators can be calculated from the Heisenberg equation.
	

\hrule
}


\newpage


{\color{violet}
\hrule
	\fbox{Derivation Multi-Mode Quantization}\\
	Strategy:\\
	1) Solve Maxwell's eqs. for the vector potential $\vec{A}$\\
	2) Write a general solution that respects all wave vectors and polarities\\
	3) Calculate E- and B-field from $\vec{A}$\\
	4) Calculate the Hamilton function\\
	5) Recognize terms in the Hamiltonian as position and momoentum\\
	5) Replace position and momentum with their respective operators\\
	6) Express the operators through the ladder operators $\hat{\alpha}_{\vec{k},s}^{(\dagger)}$\\
	\\
	Consider the two following equations known from electrodynamics:
		\begin{align*}
			-\Delta \phi - \partial_t\vec{\nabla}\vec{A} &= 0 \\
			\lr{ \frac{1}{c^2} \partial_t^2 - \Delta } \vec{A} + \vec{\nabla} \lr{ \frac{1}{c^2}\partial_t \phi + \vec{\nabla}\vec{A} } &= 0
		\end{align*}
	
	The first one follows from $\vec{\nabla}\cdot\vec{E}=0$ and the definition of the vector potentials $\vec{E} := -\vec{\nabla}\phi - \partial_t\vec{A}$. The second term arises from $\vec{\nabla}\times\vec{\nabla}\times\vec{A}$, the forth Maxwell equation and the definition of potentials.\\
	\\
	Because of the Coulomb gauge, we are free to choose
		\begin{align*}
			\phi &= 0 \\
			\vec{\nabla} \cdot \vec{A} &= 0,
		\end{align*}
	
	which simplifies the equations from before to
		\begin{align*}
			\phi &= 0 \\
			\lr{ \frac{1}{c^2} \partial_t^2 - \Delta } \vec{A} &= 0
		\end{align*}

	The latter is the wave equation for $\vec{A}$ and has the known solutions
		\begin{align*}
			\vec{A} &\propto e^{i\vec{k}_{m_x,m_y,m_z}\vec{x}}
		\end{align*}

	where $k_i = \frac{2\pi}{L} \cdot m_i$ with $m_i \in \mathbb{Z}$.\\
	\\
	The general solution then sums over all possible $\vec{k}_{m_x,m_y,m_z} \equiv \vec{k}$ and polarisations $s$.
		\begin{align*}
			\vec{A}(\vec{x},t) &= \sum_{\vec{k}} \sum_{s=1}^2 \vec{e}_{\vec{k},s} \vec{e}_{\vec{k},s} \lr{ A_{\vec{k},s}(t) \cdot e^{i\vec{k}\vec{x}} + A_{\vec{k},s}^\star(t) \cdot e^{-i\vec{k}\vec{x}} }
		\end{align*}

	Insert $\alpha$ into the equations that connect it to the B- and E-field:
		\begin{align*}
			\bullet \vec{E}(\vec{x},t) &= -\partial_t \vec{A} = \\
			&= -\partial_t \lreck{ \sum_{\vec{k}} \sum_{s=1}^2 \vec{e}_{\vec{k},s} \vec{e}_{\vec{k},s} \lr{ A_{\vec{k},s}(t) \cdot e^{i\vec{k}\vec{x}} + A_{\vec{k},s}^\star(t) \cdot e^{-i\vec{k}\vec{x}} } } = \\
			&= +i \sum_{\vec{k}} \sum_{s=1}^2 \vec{e}_{\vec{k},s} \vec{e}_{\vec{k},s} \lr{ A_{\vec{k},s}(t) \cdot e^{i\vec{k}\vec{x}} + A_{\vec{k},s}^\star(t) \cdot e^{-i\vec{k}\vec{x}} } \qquad \blacksquare \\
			\bullet \vec{B}(\vec{x},t) &= \vec{\nabla} \times \vec{A} = \vec{\nabla} \times \lreck{ \sum_{\vec{k}} \sum_{s=1}^2 \vec{e}_{\vec{k},s} \vec{e}_{\vec{k},s} \lr{ A_{\vec{k},s}(t) \cdot e^{i\vec{k}\vec{x}} + A_{\vec{k},s}^\star(t) \cdot e^{-i\vec{k}\vec{x}} } } = \\
			&{\color{red}= ... = } \\
			&= \frac{i}{c} \sum_{\vec{k},s} \omega_k \lr{ \frac{\vec{k}}{|\vec{k}|} \times \vec{e}_{\vec{k},s} } \lr{ A_{\vec{k},s} e^{ i(\vec{k}\vec{r} - \omega_k t) } - A_{\vec{k},s}^\star e^{ -i(\vec{k}\vec{r} - \omega_k t) } } \qquad \blacksquare
		\end{align*}

	The Hamilton function is then
		\begin{align*}
			H &= \int_V \diff V \, \lr{ \frac{\varepsilon_0}{2} \vec{E}^2(\vec{x},t) + \frac{1}{2\mu_0} \vec{B}^2(\vec{x},t) } {\color{red}= ... =} \sum_{\vec{k},s} \omega_k^2 A_{\vec{k},s} A_{\vec{k},s}^\star.
		\end{align*}

	Redefining the amplitudes as
		\begin{align*}
			A_{\vec{k},s} &:= \frac{1}{2\omega_k \sqrt{\varepsilon_0 V}} (\omega_k q_{\vec{k},s} + ip_{\vec{k},s} ) \\
			A_{\vec{k},s}^\star &= \frac{1}{2\omega_k \sqrt{\varepsilon_0 V}} (\omega_k q_{\vec{k},s} - ip_{\vec{k},s} ) \\
		\end{align*}
	
	allows us to recognise the Hamilton function of the harmonic oscillator.\\
	\\
	Once again, use the ladder operators to express the position and momentum operators:
		\begin{align*}
			\hat{a}_{\vec{k},s} &:= \frac{1}{ \sqrt{2\hbar\omega_{\vec{k}}} } \lr{ \omega_{\vec{k}} q_{\vec{k},s} + i p_{\vec{k},s} } \\
			\hat{a}^\dagger_{\vec{k},s} &= \frac{1}{ \sqrt{2\hbar\omega_{\vec{k}}} } \lr{ \omega_{\vec{k}} q_{\vec{k},s} - i p_{\vec{k},s} }
		\end{align*}

	\textcolor{red}{ERGÄNZEN Einsetzen}\\
\hrule
}


\newpage


\subsubsection{States of the Quantized Electromagnetic Field}
\begin{itemize}
	\item \textbf{Fock state}\\
	Fock states are the eigenstates $\ket{n}$ or $\ket{n_1, n_2, ...}$ of the single- or multi-mode Hamiltonian, thus fulfilling
		\begin{align}
			\hat{H}_n \ket{n} &= \hbar\omega_n \lr{ n + \frac{1}{2} } \ket{n} \\
			\hat{H}\ket{n_1, n_2, ...} &= \sum_{\vec{k},s} \hbar\omega_{\vec{k}} \lr{ n_{\vec{k},s} + \frac{1}{2} } \ket{n_1, n_2, ...},
		\end{align}


	\item \textbf{Thermal state}\\
	The thermal state is a mixed state with the Boltzmann statistic describing the classical probabilities.\\
	It describes a state coupled to a thermal reservoir.
		\begin{align}
			\hat{\rho}_\text{th,sm} &= \frac{1}{Z} e^{ -\frac{\hat{H}}{k_B T} } \\
			\hat{\rho}_\text{th,mm} &= \prod_{i=1}^N \frac{1}{Z}_i e^{ -\frac{\hat{H}_i}{k_B T} }.
		\end{align}
	
	
	\item \textbf{Coherent state}\\
	A coherent state $\ket{\alpha}$ is an eigenstate of the annihilation operator $\hat{a}$ with the eigenvalue $\alpha$.\\
	They are motivated by the search for states with non-vanishing E-fields at high photon occupations.\\
	Due to their constant minmal variance in the quadrature operators, they can also be seen as displaced vacuum states.
		\begin{align}
			\hat{a} \ket{\alpha} &= \alpha \ket{\alpha} \\
			\ket{\alpha} &= e^{-\frac{1}{2}|\alpha|^2} \sum_n \frac{\alpha^n}{\sqrt{n!}} \ket{n} \\
			\ket{\alpha} &= \hat{D}(\alpha) \ket{0}
		\end{align}
	
	with $\hat{D}(\alpha) := e^{\alpha \hat{a}^\dagger - \alpha^\star \hat{a}}$.


	\item \textbf{Squeezed state}\\
	Squeezed states, where the variance of one quadrature is reduced below the variance of the vacuum state at the expense of the variance of the other quadrature.
		\begin{equation}
			\ket{\xi} = \hat{S}(\xi) \ket{0},
		\end{equation}
	
	with the \textbf{squeezing operator} $\hat{S}(\xi) := e^{\frac{\xi^\star}{2} \hat{a}^2 - \frac{\xi}{2} (\hat{a}^\dagger)^2}$.


	\item \textbf{Squeezed coherent states}\\
	A squeezed coherent state is a displaced squeezed state
		\begin{equation}
			\ket{\alpha,\xi} = \hat{D}(\alpha) \hat{S}(\xi) \ket{0}.
		\end{equation}
\end{itemize}


\newpage


\begin{landscape}
\begin{table}[h]
	\centering
	\caption{Properties of single-mode Fock-, single-mode thermal, coherent, squeezed states and squeezed coherent states}
	\begin{tabular}{p{1.5cm}|c|c|c|c}
		\hline\hline
			State & $\braket{\hat{X}_1}$ & $\braket{\hat{X}_2}$ & $\Delta\hat{X}_1$ & $\Delta\hat{X}_2$ \\
		\hline\hline
			Fock & $0$ & $0$ & $\frac{1}{2}\sqrt{2n+1}$ & $\frac{1}{2}\sqrt{2n+1}$ \\
		\hline
			Thermal & $0$ & $0$ & $\frac{1}{2}\sqrt{2\braket{\hat{n}}+1}$ & $\frac{1}{2}\sqrt{2\braket{\hat{n}}+1}$ \\
		\hline
			Coherent & $\Re\lrschweif{\alpha}$ & $\Im\lrschweif{\alpha}$ & $1/2$ & $1/2$ \\
		\hline
			Squeezed & $0$ & $0$ & $\frac{1}{2} \sqrt{  e^{-2|\xi|} \cos^2( \frac{\phi}{2} ) + e^{+2|\xi|} \sin^2( \frac{\phi}{2} ) }$ & $\frac{1}{2} \sqrt{ e^{+2|\xi|} \cos^2( \frac{\phi}{2} ) + e^{-2|\xi|} \sin^2( \frac{\phi}{2} ) }$ \\
		\hline
			Squeezed Coherent & $\Re\lrschweif{\alpha}$ & $\Im\lrschweif{\alpha}$ & -''- & -''- \\
		\hline\hline
		\\
		\hline\hline
			State & $\braket{\hat{n}}$ & $\Delta \hat{n}$ & $\braket{\hat{E}}$ & $P(\tilde{n})$ \\
		\hline\hline
			Fock & $n$ & $0$ & $0$ & $\delta_{\tilde{n},n}$ \\
		\hline
			Thermal & $\frac{1}{ e^{\frac{\hbar\omega}{k_BT}}  - 1 }$ & $\sqrt{ \braket{\hat{n}}^2 + \braket{\hat{n}} }$ & $0$ & $\frac{1}{Z} \cdot e^{-\frac{E_n}{k_B T}}$ \\
		\hline
			Coherent & $|\alpha|^2$ & $|\alpha|$ & $2|\alpha| \sqrt{\frac{\hbar\omega}{2\varepsilon_0 V}} \sin\lr{ \omega t - \vec{k}\vec{x} - \varphi }$ & $e^{-|\alpha|^2} \cdot \frac{|\alpha|^{2\tilde{n}}}{\tilde{n}!}$ \\
		\hline
			Squeezed & $\sinh^2(|\xi|)$ & $2\braket{\hat{n}}( 1 + \braket{\hat{n}} )$ &  & $P(2\tilde{n}+1) = 0,\ P(2\tilde{n}) \neq 0$ \\
		\hline
			Squeezed Coherent & $|\alpha|^2 + \sinh^2(|\xi|)$ &  &  &  \\
		\hline\hline
	\end{tabular}
\end{table}

\end{landscape}


\newpage


\begin{figure}[H]
	\centering
	\fbox{
	\includegraphics[width=\textwidth]{Grafiken/quantum states.pdf}
	}
	\caption{Graphical overview of the quadrature operator's expectation values and standard deviations}
\end{figure}


\newpage


{\color{violet}
\hrule
	\fbox{Derivation Coherent States}\\
	zz: $\ket{\alpha} = e^{-\frac{1}{2}|\alpha|^2} \sum_n \frac{\alpha^n}{\sqrt{n!}} \ket{n}$\\
	\\
	Strategy:\\
	1) Expand $\alpha$ in the Fock states\\
	2) Find a recursive relation between the coefficients\\
	3) Find the coefficients $c_m(c_0)$ via proof of induction\\
	4) Find $c_0$ via normalisation condition\\
	\\
	The expansion of the coherent state in the fock basis is
		\begin{align*}
			\ket{\alpha} &= \sum_n c_n \ket{n}.
		\end{align*}
	
	Insertion into the defining eigenvalue function of the coherent states yields:
		\begin{align*}
			\hat{a} \ket{\alpha} &= \hat{a} \sum_n c_n \ket{n} = \sum_n c_n \sqrt{n} \ket{n-1} = \\
			= \alpha \ket{\alpha} &= \alpha \sum_n c_n \ket{n} \\
			\stackrel{\comment{\cdot \bra{m}}}{\Leftrightarrow} c_{m+1} \cdot \sqrt{m+1} &= \alpha \cdot c_m
		\end{align*}

	Consider the first few values for $m$:
		\begin{align*}
			\bullet \ c_1 &= \frac{\alpha}{\sqrt{1}} \cdot c_0 \\
			\bullet \ c_2 &= \frac{\alpha}{\sqrt{2}} \cdot c_1 = \frac{\alpha}{\sqrt{2}} \cdot \frac{\alpha}{\sqrt{1}} c_0 = \frac{\alpha^2}{\sqrt{2\cdot 1}} c_0 \\
			\bullet \ c_3 &= \frac{\alpha}{\sqrt{3}} \cdot c_2 = \frac{\alpha}{\sqrt{3}} \cdot \frac{\alpha^2}{\sqrt{2\cdot 1}}c_0 = \frac{\alpha^3}{\sqrt{3\cdot2\cdot1}}c_0
		\end{align*}
		
	Thus, assume
		\begin{align*}
			c_m &= \frac{\alpha^m}{\sqrt{m!}}c_0.
		\end{align*}
	
	Proof this via induction.\\
	IA: $m=1$
		\begin{align*}
			c_1 &= \alpha \cdot c_0 \stackrel{!}{=} \frac{\alpha^1}{\sqrt{1!}} \cdot c_0 = \alpha \cdot c_0 \qquad \checkmark
		\end{align*}
	
	IS: $c_{m+1} = \frac{\alpha}{\sqrt{m+1}} c_m \stackrel{!}{=} \frac{\alpha^{m+1}}{\sqrt{(m+1)!}} \cdot c_0$	
		\begin{align*}
			c_{m+1} &= \frac{\alpha}{\sqrt{m+1}} \cdot c_m = \frac{\alpha}{\sqrt{m+1}} \cdot \frac{\alpha^m}{\sqrt{m!}} \cdot c_0 = \frac{\alpha^{m+1}}{\sqrt{(m+1)!}} \cdot c_0 \qquad \blacksquare
		\end{align*}

	\textcolor{red}{ERGÄNZEN}
\hrule
}


\newpage


\subsubsection{Coherence}
The \textbf{$m$-th order coherence} function is defined as
	\begin{equation}
		G^{(m)}(x_1, ..., x_m; x_1', ..., x_m') := \text{Sp} \lrschweif{ \hat{\rho} \hat{\vec{E}}^-(x_1) ... \hat{\vec{E}}^-(x_m) \hat{\vec{E}}^+(x_1') ... \hat{\vec{E}}^+(x_m') },
	\end{equation}

where $x_i = (\vec{x}_i,t_i)$ and $x_i = (\vec{x}_i',t_i')$.\\
\\
Two physical applications are:
	\begin{itemize}
		\item \textbf{capability of interference}\\
		$|g^{(1)}|$ is equivalent to the \textbf{visibility} $V$ of an interference pattern.
			\begin{align}
				G^{(1)}(x_1,x_1') &= \text{Sp} \lrschweif{ \hat{\rho} \hat{\vec{E}}^-(x_1) \hat{\vec{E}}^+(x_1') } \\
				|g^{(1)}(x_1,x_1')| &= \frac{ G^{(1)}(x_1,x_1') }{ \sqrt{ G^{(1)}(x_1,x_1) \cdot G^{(1)}(x_1',x_1') } }
			\end{align}
		
		
		\item \textbf{classicality of the state}\\
		The measurement of the $g^{(2)}$ parameter in the \textbf{Hanbury-Brown-Twiss interferometer} yields the Mandel $Q$ parameter.
			\begin{align}
				G^{(2)}(x_1,x_2,x_1',x_2') &= \text{Sp} \lrschweif{ \hat{\rho} \hat{\vec{E}}^-(x_1) \hat{\vec{E}}^-(x_2) \hat{\vec{E}}^+(x_1) \hat{\vec{E}}^+(x_2) } \\
				|g^{(2)}(x_1,x_2)| &= \frac{G^{(2)}(x_1,x_2)}{ G^{(1)}(x_1,x_1) \cdot G^{(1)}(x_2,x_2) } \\
				g^{(2)}(\tau) &= \frac{ \braket{ \hat{\vec{E}}^-(t) \hat{\vec{E}}^-(t+\tau) \hat{\vec{E}}^+(t+\tau) \hat{\vec{E}}^+(t) } }{ \braket{ \hat{\vec{E}}^-(t) \hat{\vec{E}}^+(t) } }
			\end{align}
	\end{itemize}

\textcolor{blue}{$g^{(2)}$ is a measure for single photon sources.}


\newpage


\subsubsection{Quantum Mechanical Light-Matter Interaction}
\subsubsection*{Jaynes Cummings Hamiltonian}
The simplest, fully quantum mechanical model for light-matter interaction considers a
	\begin{itemize}
		\item \textbf{two-level atom}
		
		\item \textbf{single mode} the quantized em. field
		
		\item \textbf{dipole approximation}\\
		Assumes $\vec{A}(\vec{x},t) \approx \vec{A}(t)$ due to $\vec{k} \cdot \vec{x} \ll 1$.\\
		Used in the derivation of the interaction Hamiltonian.
		
		\item \textbf{rotating wave approximation (RWA)}\\
		Negelects fast oscillating terms $\omega+\omega_0$.\\
		\textcolor{red}{Good approximation for long timescales and for small interaction energies.}\\
		Also removes non energy conserving processes ($\ket{e}\ket{n} \to \ket{g}\ket{n-1}$ and $\ket{g}\ket{n} \to \ket{e}\ket{n+1}$)


		\item \textbf{non-relativistic}
	\end{itemize}

which leads to the Jaynes-Cummings-Hamiltonian
	\begin{align}
		\hat{H}_\text{JC} &= \hat{H}_\text{atom} + \hat{H}_\text{field} + \hat{H}_\text{int},
	\end{align}

which can be written in numerous ways, using:
\begin{itemize}
	\item \textbf{•}
		\begin{equation}
			\hat{H}_\text{JC} = \frac{\hbar\omega_0}{2} ( - \ket{g}\bra{g} + \ket{e}\bra{e} )
		+
		\hbar\omega \hat{a}^\dagger\hat{a} 
		+
		\hbar\lambda \lr{ \ket{e}\bra{g} + \ket{g}\bra{e} } \lr{ \hat{a} + \hat{a}^\dagger }
		\end{equation}
	\item \textbf{Pauli matrices}\\
		\begin{equation}
			\hat{H}_\text{JC} = \frac{\hbar\omega_0}{2} \hat{\sigma}_3
			+
			\hbar\omega \hat{a}^\dagger\hat{a}
			+
			\hbar\lambda \lr{ \ket{e}\bra{g} + \ket{g}\bra{g} } \lr{ \hat{a} + \hat{a}^\dagger }
		\end{equation}
	
	
	
	\item \textbf{conserved quantities}\\
	Conserved quantities in this simplified systems are the atomic population $\hat{P}_e := \ket{e}\bra{e} + \ket{g}\bra{g}$ and number of excitation $\hat{N}_e := \hat{a}^\dagger\hat{a} + \ket{e}\bra{e}$
		\begin{equation}
			\hat{H}_\text{JC} = \lreck{ \hbar\omega_e \hat{N}_e + \hbar\lr{ \frac{\omega_0}{2} - \omega } \hat{P}_e } + \lreck{ -\hbar \Delta \ket{g}\bra{g} + \hbar\lambda \lr{ \hat{\sigma}_+ \hat{a} + \hat{\sigma}_- \hat{a}^\dagger } }
		\end{equation}
	
	
	\item \textbf{sum of subspaces}\\
	Since the Jaynes-Cummings-Hamiltonian couples the bare states $\{\ket{g,n}, \ket{e,n-1}\}$ for arbitrary excitations of the mode $n$, one can express the Hamiltonian the following way
		\begin{align}
			\hat{H}_\text{JC} &= \sum_n \hat{H}_n = \\
			&= \sum_n \lrschweif{ \hbar\omega\lr{ n + \frac{1}{2} }
		\begin{pmatrix}
			1 & 0 \\
			0 & 1
		\end{pmatrix}
		+
		\frac{\hbar}{2}
		\begin{pmatrix}
			-\Delta & 2\lambda\sqrt{ n+1 } \\
			2\lambda\sqrt{ n+1 } & \Delta
		\end{pmatrix}
		 },
		\end{align}
	
	with the detuning $\Delta := \omega - \omega_0$.
\end{itemize}



The eigenvalues and eigenstates (\textbf{dressed states}) of the $\hat{H}_n$ Hamiltonian are
	\begin{align}
		E_{\pm,n} &= \hbar\omega \lr{ n + \frac{1}{2} } \pm \frac{\hbar}{2}\Omega_n \\
		\ket{+,n} &= \cos\theta_n \ket{e,n} + \sin\theta_n \ket{g,n+1} \\
		\ket{-,n} &= -\sin\theta_n \ket{e,n} + \cos\theta_n \ket{g,n+1}
	\end{align}

with the \textbf{generalized Rabi frequency} $\Omega_n := \sqrt{ \Delta^2 + 4\lambda^2(n+1) }$ and the rotation angles $\cos\theta_n := \frac{ 2\lambda\sqrt{n+1} }{ \sqrt{ (\Omega_n - \Delta )^2 + 4\lambda^2(n+1) } }$, $\sin\theta_n := \frac{ \Omega_n - \Delta }{ \sqrt{ (\Omega_n - \Delta )^2 + 4\lambda^2(n+1) } }$.\\
\\
Two significant results are:
	\begin{itemize}
		\item \textbf{vacuum Rabi oscillations}\\
		Even if the mode has zero excitations $n=0$, the probability for the excited state oscillates
		
		
		\item \textbf{oscillation rivival}\\
		If the mode is in a superposition of Fock states, the probability for the excited state can die out and revive again after a characteristic revival time
	\end{itemize}


\begin{figure}[h]
	\centering
	\fbox{
	\includegraphics[width=\textwidth]{Grafiken/jaynes cummings rabi oscillations.png}
	}
	\caption{Plot of the revival of a Rabi oscillation}
\end{figure}


\newpage


{\color{violet}
\hrule
	\fbox{Derivation}\\
	Strategy:\\
	1) Formulate the Hamiltonian in the electromagnetic field \\
	2) Use the dipole approximation ($\vec{A} \neq \vec{A}(t)$) \\
	3) Gauge transformation and choice of gauge \\
	4) Use operator representations of E and d\\
	\\
	Consider the Hamiltonian of a particle in an electromagnetic field. Due to minimal substitution, it reads
		\begin{align*}
			\hat{H} &= \frac{1}{2m} \lreck{ \hat{\vec{p}} - q\vec{A}(\vec{x},t) }^2 + q\varphi(\vec{x},t) + V(\vec{x})
		\end{align*}
	
	Consider a \underline{charge-free} space, which implies a vanishing electric potential $\varphi=0$.\\
	The solution of the vector potential are plane waves $\vec{A}(\vec{x},t) = \vec{A}_0 e^{i(\vec{k}\vec{x} - \omega t)} + c.c.$\\
	\\
	This solution can be approximated by realizing that $\vec{k}\vec{x} \ll 1$, which is called the \textbf{dipole approximation}, yielding $\vec{A}(\vec{x},t) \approx \vec{A}(t) = e^{i\omega t} + c.c.$\\
	\\
	It can be shown, that the Hamiltonian is invariant under the Gauge transformations
		\begin{alignat*}{2}
			\varphi(\vec{x},t) & \to \varphi'(\vec{x},t) &&= \varphi(\vec{x},t) - \partial_t \varphi(\vec{x},t) \\
			\vec{A}(\vec{x},t) & \to \vec{A}'(\vec{x},t) &&= \vec{A}(\vec{x},t) + \vec{\nabla} \chi(\vec{x},t)
		\end{alignat*}

	Choose the gauge $\chi(\vec{x},t) = -\vec{A}(t) \cdot \vec{x}$.\\
	Thus, $\vec{\nabla} \chi(\vec{x},t) = -\vec{A}(t)$.\\
	The transformed Hamiltonian then becomes:
		\begin{align*}
			\hat{H}' &= \frac{1}{2m} \lreck{ \hat{\vec{p}} - q\vec{A}'(t) }^2 + q \varphi'(\vec{x},t) + V(\vec{x}) = \\
			&= \frac{1}{2m} \lreck{ \hat{\vec{p}} - q \lr{ \vec{A}(\vec{x},t) + \vec{\nabla} \chi(\vec{x},t) } }^2 + q \lr{ \varphi(\vec{x},t) - \partial_t \chi(\vec{x},t) } + V(\vec{x}) = \\
			&\stackrel{\comment{\chi}}{=} \frac{1}{2m} \lreck{ \hat{\vec{p}} - q \lr{ \cancel{ \vec{A}(\vec{x},t) } - \cancel{ \vec{A}(t) } } }^2 + q \lr{ \cancel{ \varphi(\vec{x},t) } - \partial_t (-\vec{x} \cdot \vec{A}(t)) } + V(\vec{x}) = \\
			&= \frac{1}{2m} \hat{\vec{p}}^2 + q \cdot \vec{x} \cdot \partial_t \vec{A}(t) + V(\vec{x}) =
			\intertext{Use the definition of the electromagnetic potentials $\vec{E} = -\vec{\varphi} - \partial_t \vec{A} = -\partial_t \vec{A}$.}
			&= \frac{1}{2m} \hat{\vec{p}}^2 - q \cdot \vec{x} \cdot \vec{E} + V(\vec{x}) \equiv \\
			&\equiv \hat{H}_0 - \vec{d} \cdot \vec{E}(t),
		\end{align*}
	
	where $\vec{d} := q \cdot \vec{x}$. $\qquad \blacksquare$
\hrule
}


\newpage


\subsubsection*{Single Mode with Reservoir}
Consider a system a single mode that is coupled to a reservoir
	\begin{equation}
		\hat{H} = \hat{H}_S + \hat{H}_R + \hat{H}_\text{int} = \hbar\omega \hat{a}^\dagger\hat{a} + \sum_j \hbar\omega_j \hat{b}_j^\dagger \hat{b}_j + \sum_j \hbar\lambda_j \lr{ \hat{a}^\dagger\hat{b}_j + \hat{b}_j^\dagger\hat{a} },
	\end{equation}

where the states of the system and the reservoir are described by $\hat{\rho}_{RS}$.

	\begin{itemize}
		\item \textbf{no correlation}\\
		At $t=0$ system and reservoir state can be regarded as uncorrelated
			$$\hat{\rho}_{RS}(t) = \hat{\rho}_R(t) \otimes \hat{\rho}_S(t)$$
		
		
		\item \textbf{thermal state}\\
		The density matrix for the reservoir is assumed to be described by a thermal state
		
		
		\item \textbf{Born approximation}\\
		The density matrix for the reservoir is assumed to be time independent
		
		
		\item \textbf{continuum transition}\\
		Consider the frequencies of the reservoir as a continuum
	\end{itemize}

This leads to the \textbf{master equation}
	\begin{equation}
		\frac{\diff \hat{\rho}_S,I}{\diff t} = - \frac{\gamma}{2} \lr{ \hat{\rho}_{S,I}(t) \hat{a}^\dagger \hat{a} + \hat{a}\hat{a}^\dagger \hat{\rho}_{S,I}(t) - 2 \hat{a} \hat{\rho}_{S,I}(t) \hat{a}^\dagger }.
	\end{equation}

Important results for a single omde coupled to a reservoir are
	\begin{align}
		\lim\limits_{t\to\infty} \braket{ \hat{a}_\omega^\dagger \hat{a}_\omega }(t) &= \braket{ \hat{b}^\dagger \hat{b} }(\omega) \\
		\lim\limits_{t\to\infty} \Delta\hat{X}_1 \Delta\hat{X}_2 &= \frac{\hbar}{2} \lr{ 2\braket{\hat{n}} + 1 }
	\end{align}

\vspace{0.5cm}

{\color{violet}
\hrule
	\fbox{Derivation}\\
	Strategy:\\
	1) Consider $\hat{\rho}_{RS}$ in the interaction picture \\
	2) Calculate the Heisenberg equation\\
	3) Integrate the Heisenberg equation\\
	4) Make a few approximations\\
	\\
	\textcolor{red}{ERGÄNZEN}
\hrule
}


%\newpage
%
%
%\subsubsection{Quantum Mechanical Laser Oscillations}
%\textcolor{red}{ERGÄNZEN}




\newpage


\section{Linewidth}
The cross section is an important property because it is directly tied to the gain coefficient of the medium as well as the Einstein coefficients.
	\begin{itemize}
		\item \textbf{cross section} $\sigma(\nu)$
		
		
		\item \textbf{oscillator strength} $S$\\
		Area beneath the cross section
			\begin{equation}
				S := \int_0^\infty \sigma(\nu) \, \diff \nu
			\end{equation}
		

		\item \textbf{line shape function} $l(\nu)$\\
		Normalized cross section
			\begin{equation}
				l(\nu) := \frac{\sigma(\nu)}{S}
			\end{equation}

		
		\item \textbf{linewidth} $\Delta\nu$\\
		Full width half maximum (FWHM) of the lineshape function
			\begin{equation}
				\Delta\nu := \text{FWHM} \lreck{ l(\nu) }
			\end{equation}
	\end{itemize}


\subsection{Phenomenological Connection to Einstein Coefficients}
For the emission/absorption into/from a single mode, the relation between the rates $p_i$ and the cross section $\sigma(\nu)$ are derived phenomenologically:
	\begin{align}
		p_A &= \frac{c}{V} \sigma(\nu) \\
		p_{B_{21}} &= n \cdot \frac{c}{V} \sigma(\nu) \\
		p_{B_{21}} &= p_{B_{12}},
	\end{align}

where $n$ indicates the number of photons already present in the mode.\\
\\
For the emission/absorption into/from any mode, one derives from the previous equations:
	\begin{align*}
		A &= \int_0^\infty \diff\nu \, \lr{ \frac{c}{V} \sigma(\nu) } \cdot \lr{ V \cdot M(\nu) } \\
		B_{21} u(\nu) &= \int_0^\infty \diff\nu \, \lr{ n \cdot \frac{c}{V} \sigma(\nu) } \cdot \lr{ V \cdot M(\nu) }\\
		B_{12} u(\nu) &= B_{21} u(\nu)
	\end{align*}


\newpage


\subsection{Linewidth Broadening}
Line width broadening is the phenomenon that emission or absorption lines widen, indicating that a range of frequencies are involved in the process, not just a single one.\\
\\
Causes for line width broadening are:
	\begin{itemize}
		\item \textbf{lifetime broadening}\\
		Fundamental broadening $\Delta\nu$ of spectral lines caused by the finite lifetime of an excited state
			\begin{equation}
				\Delta E \cdot \Delta t \geq h
			\end{equation}

		{\color{violet}The Einstein coefficient connects life time and cross section.}


		\item \textbf{collisional broadening}\\
		Broadening mechanism due to collisions between atoms or molecules


		\item \textbf{Doppler broadening}\\
		The Doppler effect shifts the frequencies of emitted light, thus widening the linewidth
			\begin{equation}
				\bar{l}(\nu) = \int_{-\infty}^{+\infty} l\lr{ \nu - \nu_0\frac{v}{c} } p(v) \, \diff v
			\end{equation}
			
		
		\item \textbf{variation of dopants} (in semiconductor lasers)
	\end{itemize}

These causes can be classified into \textbf{homogeneous} and \textbf{inhomogeneous broadening}. The first subsumes broadenings, that affects all atom types equally and the latter all broadenings that differ from atom type to atom type.\\
\\
Homogeneous broadenings are described by the \textbf{Lorentz profile}
	\begin{equation}
		L(\nu; \gamma) = \frac{\gamma}{\pi \lr{ \nu^2 + \gamma^2 }}
	\end{equation}

whereas inhomogeneous broadenings are described by \textbf{Gaussian profiles}
	\begin{equation}
		G(\nu; \sigma, \nu_0) = \frac{1}{\sqrt{2\pi} \sigma} e^{- \frac{ (\nu - \nu_0)^2 }{ 2\sigma^2 } }.
	\end{equation}

If both broadenings occur, the lineshape is described by a \textbf{Voigt profile}
	\begin{equation}
		V(\nu; \gamma, \sigma) = \int_{-\infty}^{+\infty} \diff \tau \, G(\tau) \cdot L(\nu - \tau).
	\end{equation}


\newpage


\section{Resonator Optics}
A resonator is an instrument that can store light at specific frequencies which are defined by the condition for positive feedback for the phase $\phi$
	\begin{equation}
		\phi(z=2d) \stackrel{!}{=} \phi(z=0) + 2\pi \cdot n,
	\end{equation}

where $\phi(z=2d)$ can accumulate additional phases due to e.g. mirror reflections and $n\in\mathbb{N}$.\\
\\
Important properties of the resonator are:
	\begin{itemize}
		\item \textbf{resonance frequencies}\\
		Resonance frequencies are the specific frequencies of a resonator $\nu_n$.\\
		\textcolor{blue}{Example Fabry-Pérot: $\nu_n = \frac{c}{2d}, n\in\mathbb{N}$}

	
		\item \textbf{free spectral range (FSR)}\\
		Distance between the modes in frequency space $\nu_F$.\\
		\textcolor{blue}{Example Fabry-Pérot: $\nu_F = \frac{c}{2d}$}


		\item \textbf{spectral width of the modes} $\delta\nu$


		\item \textbf{finesse}\\
		Measure of the maximum number of roundtrips unamplified light can take inside a resonator:
			\begin{equation}
				\mathcal{F} := \frac{\nu_F}{ \delta\nu }
			\end{equation}
		
		with the free spectral range $\nu_F$ and the spectral width of the modes $\delta\nu$.


		\item \textbf{quality factor}\\
		The quality factor is proportional to the storage time of the photons in units of the optical period
			\begin{equation}
				Q := 2\pi \frac{\text{stored energy}}{\text{energy loss per period}} {\color{red}=} \frac{\nu_0}{\nu_F} \mathcal{F}
			\end{equation}


		\item \textbf{intensity}\\
		\textcolor{blue}{Example Fabry-Pérot: $I = \frac{I_0}{1 - 2|r|\cos\phi + |r|^2} \equiv \frac{I_\text{max}}{1 + \lr{ \frac{2\mathcal{F}}{\pi} }^2 \sin^2(\phi/2)}$ with $I_\text{max} = \frac{I_0}{1 - |r|^2}$.}


		\item \textbf{mode density}\\
		Indicates the number of modes per frequency per resonator width $M(\nu)$


		\item \textbf{mode volume}\\
		Ratio between the energy of a mode devided by the maximum energy density
			\begin{equation}
				V_\text{vol} = \frac{U}{u_\text{max}}
			\end{equation}
	\end{itemize}


%\begin{figure}[H]
%	\centering
%	\includegraphics[width=0.4\textwidth]{Grafiken/resonator intensity.pdf}
%	\includegraphics[width=0.4\textwidth]{Grafiken/resonator intensity 2.pdf}
%	\caption{Intensity of a resonator for two different finesses; the left intensity spectrum has a higher finesse than the right}
%\end{figure}


\newpage


\subsection{Resonator Types}
\begin{itemize}
	\item \textbf{Fabry-Pérot resonator}\\
	Resonator consisting of two plane mirrors facing each other


	\item \textbf{spherical resonators}\\
	Resonators with curved mirror surfaces


	\item \textbf{ring resonators}\\
	Circular resonators, that do not reflect light back onto itself
	
	
	\item \textbf{bow-tie resonators}\\
	Ring resonators, that are folded into themselves
	
	
	\item \textbf{whispering gallery mode resonators}\\
	Resonator that makes use of total internal reflection to trap light at the border


	\item \textbf{distributed-feedback resonators}\\
	Resonators that create a positive feedback using periodic structures
	
	
	\item \textbf{photonic crystal resonators}\\
	Periodic structure that creates bandgaps for light, thus trapping it
\end{itemize}


\subsection{Stability Condition of Spherical Resonators}
A resonator is stable, when the light ray stays close to the optical axis after many revolutions, which is equivalent to
	\begin{equation}
		0 \leq g_1 g_2 \leq 1
	\end{equation}

where $g_i := 1 + \frac{d}{R_i}$ and $R_i$ is the radius of curvature.

\begin{figure}[H]
	\centering
	\fbox{
	\includegraphics[width=\textwidth]{Grafiken/stability condition.pdf}
	}
	\caption{Stability diagram $g_2(g_1)$ with the corresponding mirror geometries}
\end{figure}


\newpage


{\color{violet}
\hrule
	\fbox{Derivation}\\
	zz: $0\leq g_1 g_2 \leq 1$\\
	\\
	Strategy:\\
	1) Describe the round trip via ABCD formalism\\
	2) Solve for $m$ roundtrips\\
	3)Establish a condition for which the solution doesn't diverge\\
	\\
	The roundtrip light ray after one roundtrip in a spherical resonator can be described by the following matrix multiplication
		\begin{align*}
			\begin{pmatrix}
				x' \\ \theta'
			\end{pmatrix}
			&=
			\begin{pmatrix}
				1 & 0 \\
				\frac{2}{R_2} & 1
			\end{pmatrix}
			\cdot
			\begin{pmatrix}
				1 & d \\
				0 & 1
			\end{pmatrix}
			\cdot
			\begin{pmatrix}
				1 & 0 \\
				\frac{2}{R_2} & 1
			\end{pmatrix}
			\cdot
			\begin{pmatrix}
				1 & d \\
				0 & 1
			\end{pmatrix}
			\cdot
			\begin{pmatrix}
				x \\ \theta
			\end{pmatrix}
			=\\
			&= 
			\begin{pmatrix}
				1 + \frac{2d}{R_2} & d + d \lr{ 1 + \frac{2d}{R_2} } \\
				\frac{2}{R_1} + \frac{2}{R_2} \lr{ 1 + \frac{2d}{R_1} } & 1 + \frac{2d}{R_1} + d \lr{ \frac{2d}{R_1} + \frac{2}{R_2} \lr{ 1 + \frac{2d}{R_1} } }
			\end{pmatrix}
			\cdot
			\begin{pmatrix}
				x \\ \theta
			\end{pmatrix}
			\equiv \\
			&\equiv
			\begin{pmatrix}
				A & B \\
				C & D
			\end{pmatrix}
			\cdot
			\begin{pmatrix}
				x \\ \theta
			\end{pmatrix}
		\end{align*}
	
	$m$ roundtrips correspond to
		\begin{align*}
			\begin{pmatrix}
				x_{m+1} \\ \theta_{m+1}
			\end{pmatrix}
			&=
			\begin{pmatrix}
				A & B \\
				C & D
			\end{pmatrix}
			\cdot
			\begin{pmatrix}
				x_m \\ \theta_m
			\end{pmatrix} \\
			\Leftrightarrow
			&
			\begin{cases}
				x_{m+1} &= A x_m + B \theta_m \\
				\theta_{m+1} &= C x_m + D \theta_m
			\end{cases}
		\end{align*}
	
	Rearranging both equations for $\theta_m$ and setting these relations equal to each other yields:
		\begin{align*}
			\theta_{m+1} &= \frac{x_{m+1} - A y_m}{B} D + C x_m
		\end{align*}
	
	Then, use the equation $x_m+1 = A x_m + B \theta_m$, shift $m\to m+1$, rearrange for $\theta_{m+1}$ and insert it into the previous expression:
		\begin{align*}
			\frac{x_{m+2} - A x_{m+1}}{B} &= \frac{x_{m+1} - A x_m}{B}D + C x_m \\
			\Leftrightarrow x_{m+2} &= (A+D) x_{m+1} + (BC - AD) x_m
		\end{align*}
	
	Define $F := AD-BC$ and $b := \frac{A+D}{2}$:
		\begin{align*}
			x_{m+2} &= 2b x_{m+1} - F^2 x_m
		\end{align*}
	
	Insert the ansatz $x_m = x_0 h^m$:
		\begin{align*}
			\cancel{x_0} h^{m+2} &= 2b \cancel{x_0} h^{m+1} - F^2 \cancel{x_0} h^m \\
			\Leftrightarrow h^2 - 2b \cdot h + F^2 &= 0
		\end{align*}


\newpage


	The solution to this quadratic equation is
		\begin{align*}
			h_\pm &= b \pm \sqrt{F^2 - b^2} = b \pm i F \sqrt{ 1 - \lr{ \frac{b}{F} }^2 } = b \pm i F \sqrt{ 1 - \cos^2\varphi } = F\cos\varphi \pm i F \sin\varphi =  
			&= F e^{\pm i \varphi}
		\end{align*}
	
	Insert this into the solution of $x_m$:
		\begin{align*}
			x_m &= x_0 F^m e^{\pm i m \varphi}
		\end{align*}
	
	The general solution is then a linear combination, thus:
		\begin{align*}
			x_m &= x_\text{max} F^m \sin(m\varphi + \varphi_0)
		\end{align*}
		
	In our case $F=1$ \textcolor{gruen}{(see the result of the matrix multiplication in the first step)} and $b$ becomes
		\begin{align*}
			b &= \frac{A+D}{2} = 2 + \frac{4d}{R_1} + \frac{4d}{R_2} + \frac{4d^2}{R_1R_2} = ...  = 2 \lr{ 1 + \frac{d}{R_1} } \lr{ 1 + \frac{d}{R_2} } - 1
		\end{align*}
	
	$x_m$ does not diverge if the sine stays a sine, which means that the argument has to stay real. This is the case if $|b|$ doesnt exceed 1, because $b$ is the $\arccos$ function, which will return imaginary values for $|b|>1$, thus turning sine into sinh. The latter would be a divergent case.\\
	Using this condition:
		\begin{align*}
			-1 &\leq |b| \leq 1 \\
			\Leftrightarrow -1 &\leq 2 \lr{ 1 + \frac{d}{R_1} } \lr{ 1 + \frac{d}{R_2} } - 1 \leq 1 \\
			\Leftrightarrow 0 &\leq 2 \lr{ 1 + \frac{d}{R_1} } \lr{ 1 + \frac{d}{R_2} } 2 \\
			\Leftrightarrow 0 &\leq \lr{ 1 + \frac{d}{R_1} } \lr{ 1 + \frac{d}{R_2} } 1 \\
			\Leftrightarrow -1 &\leq g_1 g_2 \leq 1 \qquad \blacksquare
		\end{align*}
\hrule
}


\newpage


\subsection{Resonator Losses}
The loss in a resonator roundtrip arises from imperfect mirrors $R_1$, $R_2$ and absorption through the medium described by the absorption coefficient $\alpha_s$
	\begin{equation}
		|r|^2 = R_1 \cdot R_2 \cdot e^{-2\alpha_s d} \equiv e^{-2\alpha_R d},
	\end{equation}

where $d$ is the length of the resonator and $\alpha_R := \alpha_s + \alpha_{R_1} + \alpha_{R_2}$.\\
\\
In this context, one makes the definition of the photon lifetime
	\begin{equation}
		\tau_P := \frac{1}{c\alpha_R},
	\end{equation}

which is the decaying constant for $I(t)$.






\newpage


\subsection{Gaussian Beam}
The Gaussian beam is a solution to the paraxial Helmholtz equation and occur in spherical resonators and can be expressed as

	\begin{align}
		U(\rho,z) &= \frac{A(z)}{q(z)} \cdot \exp \lrschweif{ -ik \frac{\rho^2}{2q(z)} } \cdot \exp\lrschweif{ -ikz } \\
		U(\rho,z) &= A_0 \frac{W_0}{W(z)} \cdot \exp \lrschweif{ -\frac{\rho^2}{W^2(z)} } \cdot \exp \lrschweif{ -ikz - ik\frac{\rho^2}{2R(z)} + i\zeta(z) },
	\end{align}
	
with the terms $q(z) := z + i z_0$,
	\begin{itemize}
		\item \textbf{beam radius}
			\begin{equation}
				W(z) = W_0 \sqrt{ 1 + \lr{ \frac{z}{z_0} }^2 }
			\end{equation}


		\item \textbf{waist/waist radius}
			\begin{equation}
				W_0 = \sqrt{ \frac{ \lambda z_0 }{\pi} }
			\end{equation}


		\item \textbf{distance from the optical axis}
			\begin{equation}
				\rho = \sqrt{ x^2 + y^2 }
			\end{equation}


		\item \textbf{curvature}
			\begin{equation}
				R(z) = z \lreck{ 1 + \lr{\frac{z_0}{z}}^2}
			\end{equation}


		\item \textbf{Gouy phase}
			\begin{equation}
				\zeta(z) = \arctan\lr{\frac{z}{z_0}}
			\end{equation}
	\end{itemize}


\begin{figure}[H]
	\centering
	\fbox{
	\includegraphics[width=0.4\textwidth]{Grafiken/gaussian beam electric field.pdf}
	\includegraphics[width=0.4\textwidth]{Grafiken/gaussian beam intensity.pdf}
	}
	\caption{Qualitative plots of the electric field $\Re\lrschweif{U}$ and intensity $I = |U|^2$ of a gaussian beam}
\end{figure}


\newpage


The Gaussian beam has the following properties:
	\begin{itemize}
		\item \textbf{intensity}\\
		The intensity for a fixed $z$ follows a Gaussian distribution is
			\begin{equation}
				I(\rho,z) \propto \exp \lrschweif{-2 \frac{\rho^2}{W^2(z)} }
			\end{equation}


		\item \textbf{half-opening angle}\\
			\begin{equation}
				\theta_0 = \sqrt{ \frac{\lambda}{\pi z_0} }
			\end{equation}


		\item \textbf{focus length/confocal parameter}\\
		Axial distance within which the beam radius is smaller than $\sqrt{2}W_0$
			\begin{equation}
				z_0 = \frac{\pi W^2}{\lambda}
			\end{equation}
		
		\textcolor{blue}{Beispiel: $\lambda = 633 \, \text{nm}$, $2W_0 = 2 \, \text{cm}$ $\Rightarrow 2z_0= 1km$}


		\item \textbf{curvature}\\
		- The wavefront has curvature $R\to\infty$ at $z=0$ and $z\to\pm\infty$\\
		- The wavefront is most curved at the focal length
			\begin{equation}
				R(z=z_0) = 2 z_0.
			\end{equation}
	\end{itemize}

\begin{figure}[H]
	\centering
	\includegraphics[width=\textwidth]{Grafiken/gaussian beam.jpg}
	\caption{Scheme properties of the Gaussian beam}
\end{figure}



\newpage


Properties concerning a Gaussian beam inside a spherical resonator are:
	\begin{itemize}
		\item \textbf{''reflection points''}\\
		At the positions of the mirrors $z_1$, $z_2$, the curvature of the Gaussian matches with the mirror curvatures $R(z_1) = R_1$, $R(z_2) = R_2$
		
		
		\item \textbf{resonance frequency}
			\begin{equation}
				\nu_n = \lr{ n + \frac{\Delta\zeta}{\pi} } \cdot \frac{c}{2d},
			\end{equation}
		
		where $n\in\mathbb{N}$.
		

		\item \textbf{free spectral range}
			\begin{equation}
				\nu_F = \frac{c}{2d}
			\end{equation}
	\end{itemize}


\subsubsection*{Beam Transformation}
The transformation of a Gaussian beam due to an optical system described by the ABCD formalism changes the parameter $q(z)$ in the following way:
	\begin{equation}
		q(z) \to q'(z) = \frac{A \cdot q(z) + B}{C \cdot q(z) + D}
	\end{equation}


\newpage


\section{Pumping}
Pumping is the mechanism of exciting the electrons in the doped atoms to create a population inversion, thus providing the energy for the amplification process. There, multiple states apart from the two-states can be involved.


\subsection{Pump Types}
	\begin{itemize}
		\item \textbf{electrical discharge (DC or AC)}
		
		
		\item \textbf{flash lamp}
		
		
		\item \textbf{optical pumping}
	\end{itemize}


\subsection{Laser Rate Equations \& Pumping Schemes}
The rate equations for the populations involved in the spontaneous emission due to pumping are called \textbf{laser rate equations}.\\
\\
Different pumping schemes are
	\begin{itemize}
		\item \textbf{two-level pumping} (without light)
			\begin{equation}
				\Delta n_0 := \Delta n = R_1\tau_1 + R_2\tau_2 \lr{ 1 - \frac{\tau_1}{\tau_{21}} }
			\end{equation}
		
		with the pumping rate filling state 2 $R_2$ and depleting state 1 $R_1$ as well as the life times $\tau_i$.
		
		
		\item \textbf{two-state laser} (with light)\\
			\begin{equation}
				\Delta n = \frac{ \Delta n_0 }{ 1 + \tau_s W_\text{ind} },
			\end{equation}
		
		with the rate for stimulated emission $W_\text{ind}$, $\tau_s := \tau_1 + \tau_2 - \frac{\tau_1\tau_2}{\tau_{21}}$
		
		
		\item \textbf{three-state laser}\\
			\begin{equation}
				\Delta n \approx \frac{ \lr{ \frac{n_{at}(\tau_2 W - 1)}{ 1 + \tau_2 W } } }{ 1 + \lr{ \frac{2\tau_2}{ 1 + \tau_2 W } } W_\text{ind} },
			\end{equation}
		
		
		
		\item \textbf{four-state laser}\\
		Doesn't need pumping for a population inversion to exist:
			\begin{equation}
				\Delta n \approx \frac{ \lr{ \frac{\tau_2 n_{at} W }{ 1 + \tau_2 W } } }{ 1 + \lr{ \frac{\tau_2}{ 1 + \tau_2 W } } W_\text{ind} },
			\end{equation}
		
		with $n_{at} := n_0 + n_1 + n_2 + n_3$, the pump rate $W$.
	\end{itemize}



\newpage


\begin{figure}[h]
	\centering
	\fbox{
	\includegraphics[width=\textwidth]{Grafiken/pump scheme.pdf}
	}
	\caption{Pump scheme of a four- and three-state laser}
\end{figure}


\newpage


{\color{violet}
\hrule
	\fbox{Derivation: Empty Pumped Resonator}\\
	\begin{figure}[h]
		\centering
		\fbox{
		\includegraphics[width=0.4\textwidth]{Grafiken/rate equations.png}
		}
		\caption{Scheme for the processes in a two-state system}
	\end{figure}
	
	The rate equations for the states 1 and 2 are described by:
		\begin{align*}
			\frac{\diff N_1}{\diff t} &= -\frac{N_2}{\tau_{21}} - \frac{N_1}{\tau_1} - R_1 \\
			\frac{\diff N_2}{\diff t} &= -\frac{N_2}{\tau_2} + R_2
		\end{align*}
	
	In the stationary state $\frac{\diff N_1}{\diff t} = \frac{\diff N_2}{\diff t} = 0$, the equations become
		\begin{align*}
			0 &= -\frac{N_2}{\tau_{21}} - \frac{N_1}{\tau_1} - R_1 \\
			0 &= -\frac{N_2}{\tau_2} + R_2.
		\end{align*}
	
	The second equation yields $N_2 = R_2 \tau_2$.\\
	Inserting into the first equation gives
		\begin{align*}
			0 &= -\frac{R_2\tau_2}{\tau_{21}} - \frac{N_1}{\tau_1} - R_1 \\
			\Leftrightarrow N_1 &= -\tau_1 \lr{ \frac{R_2\tau_2}{\tau_{21}} + R_1 }.
		\end{align*}
	
	Thus, the population inversion is
		\begin{align*}
			\Delta N &= N_2 - N_1 = R_2\tau_2 + \tau_1 \lr{ \frac{R_2\tau_2}{\tau_{21}} - R_1 } = \\
			&= R_2\tau_2 \lr{ 1 - \frac{\tau_1}{\tau_{21}} } + R_1\tau_1
		\end{align*}
\hrule
}


\newpage


\section{Laser Oscillations \& Dynamics}
A laser is lasing at a frequency $\nu$, when two conditions are being met:
	\begin{itemize}
		\item \textbf{threshold condition}\\
		The gain must be greater or equal to the loss.\\
		This can also be formulated in terms of population inversion:
			\begin{align*}
				g(\nu) &\stackrel{!}{\geq} \alpha \\
				\Leftrightarrow \Delta n &\stackrel{!}{\geq} \Delta n_\text{thr} := \frac{\alpha}{\sigma(\nu)}
			\end{align*}
		
		
		\item \textbf{phase condition}\\
		The resonator geometrie leads to a phase shift in the electric field of $2\pi \cdot n$ after one roundtrip.
	\end{itemize}


\subsection{Single Mode Lasing}
\begin{figure}[h]
	\centering
	\fbox{
	\includegraphics[width=\textwidth]{Grafiken/single mode lasing.png}
	}
	\caption{Scheme of hole burning}
\end{figure}


\newpage


\subsection{Hole Burning}
\begin{figure}[h]
	\centering
	\fbox{
	\includegraphics[width=0.5\textwidth]{Grafiken/hole burning.png}
	}
	\caption{Scheme of hole burning}
\end{figure}


\subsection{Mode Selection}
\begin{figure}[h]
	\centering
	\fbox{
	\includegraphics[width=0.6\textwidth]{Grafiken/mode selection.png}
	}
	\caption{Scheme of mode selection}
\end{figure}


\newpage


\subsection{Gain Switching}
\begin{figure}[h]
	\centering
	\fbox{
	\includegraphics[width=\textwidth]{Grafiken/population inversion dynamics.pdf}
	}
	\caption{Qualitative scheme of the dynamics of the population inversion $\Delta n$ and photon flux $\phi$ reacting to the activation and switch-off of a pump}
\end{figure}


\newpage


In the following discussion, we focus on the properties population inversion $\Delta n$, population inversion without light in the resonator $\Delta n_0$, the threshold population inversion $\Delta n_\text{thr}$ as well as the photon flux $\phi$.
\begin{enumerate}
	\item \textbf{initial state}\\
		\begin{align*}
			\Delta n_0 &= 0 \\
			\Delta n(t) &= 0 \\
			\phi &= 0
		\end{align*}
	
	
	\item \textbf{pump is switched on}\\
	The pump with rate $W$ is being switched on, which is described by $\Delta n_0(W)$ (see pumping chapter).\\
	The rate equations lead to a rise in the population inversion, which creeps towards $\Delta n_0$
		\begin{align*}
			\Delta n_0 &= \text{konst.} > 0\\
			\Delta n(t) &> 0\\
			\phi &= 0
		\end{align*}
	
	
	\item \textbf{lasing begins}\\
	Spontaneous emission creates light in the resonator, which gets amplified due to the population inversion.\\
	Since $\Delta n(t) = \frac{1}{1 + \frac{I}{I_\text{sat}}}$, the amplification of light decreases $\Delta n(t)$.\\
	Due to light amplification, the photon flux rises.
		\begin{align*}
			\Delta n_0 &= \text{konst.} \\
			\Delta n(t)& > 0 \\
			\phi &> 0
		\end{align*}
	
	
	\item \textbf{steady state}\\
	When the population inversion reaches the threashold population inversion, no further light amplification occurs.\\
	Thus, the population as well as the photon flux inversion stays constant.
		\begin{align*}
			\Delta n_0 &= \text{konst.} \\
			\Delta n(t)&= \text{konst.} \\
			\phi &= \text{konst.}
		\end{align*}
	
	
	\item \textbf{pump switch-off}\\
	When the pump is beeing switched off, the decaying processes for state 2 in the pumping scheme dominates.\\
	Thus, the population inversion converges to its initial value again. Specifically, it falls below the threashold population inversion.\\
	From the latter it follows, that the absorption dominates the gain and the photon flux reduces as well.
\end{enumerate}


\newpage


\subsection{Q Switching}
\begin{figure}[h]
	\centering
	\fbox{
	\includegraphics[width=\textwidth]{Grafiken/q switching.pdf}
	}
	\caption{Scheme of q switching}
\end{figure}


\subsection{Mode Locking}
\textcolor{gruen}{Das Pferd von hinten aufzäumen: Auf den ''glücklichen Zufall'' hoffen; Absorber für nieder intensives Licht einbauen}


\newpage


\section{Laser Types}
\subsection{Amplification Types}
	\begin{itemize}
		\item \textbf{solid-state laser}\\
		Laser that uses a solid gain medium\\
		\textcolor{blue}{Example: Ruby laser}
		
		
		\item \textbf{gas laser}\\
		Laser that uses gases of molecules or ions as a gain medium\\
		\textcolor{blue}{Example: CO$_2$ laser}	
		
		
		\item \textbf{dye laser}\\
		Laser that uses a fluorescent molecule solution as a gain medium \\
		Have a broad tuning range
		
		
		\item \textbf{semi-conductor laser}\\
		DC pumped pn-junction with mirrors at the end\\
		Electron hole recombination creates laser radiation
	\end{itemize}


\subsection{Temporal Operation Types}
	\begin{itemize}
		\item \textbf{cw lasers}\\
		Output an uninterrupted beam of light with a constant ouput power over time
		
		
		\item \textbf{pulsed lasers}\\
		Output short intense bursts of light.\\
		Realized through \textbf{gain-switching}, \textbf{Q-switching} or \textbf{mode-locking}
	\end{itemize}


\newpage


\section*{Appendix A: Two-state atom}
The Hamiltonian of the two-state atom is
	\begin{equation}
		\hat{H} = \hbar\omega_g \ket{g}\bra{g} + \hbar\omega_e \ket{e}\bra{e} + \vec{E}(t) \cdot \vec{d} \lr{ \ket{g}\bra{e} + \ket{e}\bra{g} }
	\end{equation}
	
with the electric field $\vec{E}(t) = E_0 \cos(\omega t) \vec{e}$.\\
\\
An ansatz for the wave function is the superposition of the ground and exited state
	\begin{equation}
		\ket{\psi(t)} = C_g(t) \ket{g} + C_e(t) \ket{e}
	\end{equation}

which leads to a system of differential equations for $C_g(t), C_e(t)$.\\
The occupation probability of the ground and excited state are then
	\begin{alignat}{2}
		P_e &:= |C_e(t)|^2 &&= \lr{ \frac{\Omega_0}{\Omega} }^2 \sin^2\lr{ \frac{\Omega t}{2} } \\
		P_g &:= |C_g(t)|^2 &&= \cos^2\lr{ \frac{\Omega t}{2} } + \lr{ \frac{\Delta}{\Omega} }^2 \sin^2\lr{ \frac{\Omega t}{2} },
	\end{alignat}

with the \textbf{Rabi frequency} $\Omega_0 := \frac{\vec{d}\vec{e}}{\hbar}$, \textbf{detuning} $\Delta := \omega_0 - \omega$ and the \textbf{general Rabi frequency} $\Omega := \sqrt{ \Omega_0^2 + \Delta^2 }$.

\begin{figure}[h]
	\centering
	\includegraphics[width=0.6\textwidth]{Grafiken/two state atom.pdf}
	\caption{Occupation probabilities for the two-state atom and $\Omega_0 = 1$, $\Delta=0,5$}
\end{figure}


\newpage


{\color{violet}
\hrule
	\fbox{Derivation}\\
	Insert the ansatz $\ket{\psi(t)} = C_g(t) \ket{g} + C_e(t) \ket{e}$ into the Schrödinger equation:
		\begin{align*}
			i\hbar \partial_t \ket{\psi(t)} &= \hat{H} \ket{\psi(t)} \\
			\Leftrightarrow i\hbar \partial_t \lr{ C_g(t) \ket{g} + C_e(t) \ket{e} } &= \lr{ \hbar\omega_g \ket{g}\bra{g} + \hbar\omega_e \ket{e}\bra{e} + \vec{E}\vec{d} \lr{ \ket{g}\bra{g} + \ket{e}\bra{e} } } \lr{ C_g(t) \ket{g} + C_e(t) \ket{e} } \\
			\Leftrightarrow \lr{ i\hbar \partial_t C_g(t) } \ket{g} + \lr{ i\hbar \partial_t C_e(t) } \ket{e} &= \hbar\omega_g \ket{g} C_g(t) + \hbar\omega_e \ket{e} C_e(t) + \vec{E}\vec{d} \ket{e} C_g(t) + \vec{E}\vec{d} \ket{g} C_e(t)
		\end{align*}
	
	Insert the electric field $\vec{E}(t) = E_0 \cos(\omega t) \vec{e}$ and use the abbreviation $d_0 := \vec{d}\vec{e}$.\\
	Coefficient comparison regarding \fbox{$\ket{g}$}, \fbox{$\ket{e}$} yields:
		\begin{align*}
			i\hbar \partial_t C_g(t) &= \hbar\omega_g C_g(t) + E_0 \cos(\omega t) d_0 C_e(t) \\
			i\hbar \partial_t C_e(t) &= \hbar\omega_e C_e(t) + E_0 \cos(\omega t) d_0 C_g(t)
		\end{align*}

	To solve this system of differential equations, introduce the substitutions $c_e(t) := C_e(t) e^{i\omega_e t}, c_g(t) := C_g(t) e^{i\omega_g t}$ as well as $\omega_0 := \omega_e - \omega_g$.\\
	Rearrange those expressions to isolate $C_g(t), C_e(t)$ and insert them into both equations. Begin with $i\hbar \partial_t C_g(t)$:
		\begin{align*}
			i\hbar \partial_t C_g(t) &= i\hbar \partial_t \lr{ c_g(t) e^{-i\omega_g t} } = i\hbar \lr{ \frac{\diff}{\diff t} c_g(t) } e^{-i\omega_g t} + i\hbar(-i\omega_g) c_g(t) e^{-i\omega_g t} = \\
			&= i\hbar \lr{ \frac{\diff}{\diff t} c_g(t) } e^{-i\omega_g t} + \hbar\omega_g c_g(t) e^{-i\omega_g t} \stackrel{!}{=} \\
			&\stackrel{!}{=} \hbar\omega_g C_g(t) + E_0 \cos(\omega t) d_0 C_e(t) = \hbar\omega_g \lr{ c_g(t) e^{-i\omega_g t} } + \\
			&\phantom{=} + E_0 \cos(\omega t) d_0 \lr{ c_e(t) e^{-i\omega_e t} } = \\
			\stackrel{\comment{\cdot e^{i\omega_g t}}}{\Leftrightarrow} i\hbar \lr{ \frac{\diff}{\diff t} c_g(t) } + \cancel{ \hbar\omega_g c_g(t) } &= \cancel{ \hbar\omega_g c_g(t) } + E_0 \underbrace{ \cos(\omega t) }_{=\frac{1}{2} \lr{ e^{i\omega t} + e^{-i\omega t} } } d_0 c_e(t) \underbrace{ e^{-i(\omega_e - \omega_t)t} }_{\equiv e^{-i\omega_0 t}} \\
			\Leftrightarrow i\hbar \frac{\diff}{\diff t} c_g(t) &= \frac{E_0}{2} d_0 \lr{ e^{i(\omega - \omega_0)t} + e^{-i(\omega + \omega_0)t} }
		\end{align*}

	Analogously for $i\hbar\partial_t C_e(t)$:
		\begin{align*}
			\Rightarrow i\hbar \frac{\diff}{\diff t} c_e(t) &= \frac{E_0}{2}d_0 \lr{ e^{i(\omega + \omega_0)t} + e^{-i(\omega - \omega_0)t} }
		\end{align*}
	
	Now, use another substitution $\tilde{c}_g(t) := e^{i\Delta t/2} c_g(t)$, $\tilde{c}_e(t) := e^{-i\Delta t/2} c_e(t)$, $\Delta := \omega_0 - \omega$ and $\Omega_0 := \frac{E_0d_0}{\hbar}$ to simplify the equations. Again, rearrange for $c_e(t)$, $c_g(t)$ and insert them into the previous equations.
	

\newpage
	
	
	Begin with $i\hbar\frac{\diff}{\diff t}c_g(t)$:
		\begin{align*}
			i\hbar\frac{\diff}{\diff t}c_g(t) &= i\hbar\frac{\diff}{\diff t} \lr{ c_g(t) e^{-i\Delta t/2} } = i\hbar \lr{ \frac{\diff}{\diff t} c_g(t) } e^{-i\Delta t/2} + i\hbar(-i\Delta/2) c_g(t) e^{-i\Delta t/2} = \\
			&= i\hbar \lr{ \frac{\diff}{\diff t} c_g(t) } e^{-i\Delta t/2} + i\hbar(-i\Delta/2) c_g(t) e^{-i\Delta t/2} = \\
			&= i\hbar \lr{ \frac{\diff}{\diff t} c_g(t) } e^{-i\Delta t/2} + \hbar\Delta/2 c_g(t) e^{-i\Delta t/2} \stackrel{!}{=} \\
			&\stackrel{!}{=} \frac{E_0}{2}d_0 \lr{ \tilde{c}_e(t) e^{i\Delta t/2} } \lr{ e^{i(-\omega_0 + \omega)t} + e^{-i(\omega_0 + \omega)t} } \\
			\stackrel{\comment{\cdot e^{i\Delta t/2}}}{\Leftrightarrow} i\hbar \frac{\diff}{\diff t} \tilde{c}_g(t) + \frac{\hbar\Delta}{2} \tilde{c}_g(t) &= \frac{E_0}{2}d_0 \tilde{c}_e(t) e^{i\Delta t} \lr{ e^{-i\Delta t} + e^{-i(\omega_0 +\omega)t } } \\
			\Leftrightarrow i\hbar \frac{\diff}{\diff t} \tilde{c}_g(t) + \frac{\hbar\Delta}{2} \tilde{c}_g(t) &= \frac{E_0}{2}d_0 \tilde{c}_e(t) \lr{ 1 + e^{-i2\omega t} }
		\stackrel{\commenttext{RWA}}{\approx} \frac{E_0}{2}d_0 \tilde{c}_g(t)
		\intertext{where in the last step, the \underline{rotating wave approximation} (RWA) was made.}
			\stackrel{\comment{\cdot i/\hbar}}{\Leftrightarrow} \frac{\diff}{\diff t} \tilde{c}_g(t) &= - \lr{ -\frac{i}{\hbar} } \frac{\hbar\Delta}{2} \tilde{c}_g(t) + \lr{ -\frac{i}{\hbar} } \underbrace{ \frac{E_0}{2}d_0 }_{\equiv \Omega_0} \tilde{c}_e(t) \\
			\Leftrightarrow \frac{\diff}{\diff t} \tilde{c}_g(t) &= \lr{ -\frac{i}{2} } \lr{ -\Delta \tilde{c}_g(t) + \Omega_0 \tilde{c}_e(t) }
		\end{align*}
		
	Analogously:
		\begin{align*}
			\frac{\diff}{\diff t} \tilde{c}_e(t) &= \lr{ -\frac{i}{2} } \lr{ \Omega_0 \tilde{c}_g(t) + \Delta \tilde{c}_e(t) }
		\end{align*}
	
	The new system of equations can be formulated using matrix multiplication:
		\begin{align*}
			\begin{pmatrix}
				\frac{\diff}{\diff t} \tilde{c}_g(t)\\
				\frac{\diff}{\diff t} \tilde{c}_e(t)
			\end{pmatrix}
			&=
			\lr{ -\frac{i}{2} }
			\begin{pmatrix}
				-\Delta & \Omega_0 \\
				\Omega_0 & \Delta
			\end{pmatrix}
			\begin{pmatrix}
				\tilde{c}_g(t) \\
				\tilde{c}_e(t)
			\end{pmatrix}
		\end{align*}
	
	Compute the second derivative of the coefficients $\tilde{c}_g(t), \tilde{c}_e(t)$, because this will decouple the differential equations:
		\begin{align*}
			\begin{pmatrix}
				\frac{\diff^2}{\diff t^2} \tilde{c}_g(t)\\
				\frac{\diff^2}{\diff t^2} \tilde{c}_e(t)
			\end{pmatrix}
			&=
			\frac{\diff}{\diff t}
			\begin{pmatrix}
				\frac{\diff}{\diff t} \tilde{c}_g(t)\\
				\frac{\diff}{\diff t} \tilde{c}_e(t)
			\end{pmatrix}
			=
			\frac{\diff}{\diff t}
			\lr{ -\frac{i}{2} }
			\begin{pmatrix}
				-\Delta & \Omega_0 \\
				\Omega_0 & \Delta
			\end{pmatrix}
			\begin{pmatrix}
				\tilde{c}_g(t) \\
				\tilde{c}_e(t)
			\end{pmatrix}
			= \\
			&=
			\lr{ -\frac{i}{2} }
			\begin{pmatrix}
				-\Delta & \Omega_0 \\
				\Omega_0 & \Delta
			\end{pmatrix}
			\begin{pmatrix}
				\frac{\diff}{\diff t} \tilde{c}_g(t) \\
				\frac{\diff}{\diff t} \tilde{c}_e(t)
			\end{pmatrix}
			=
			\lr{ -\frac{i}{2} }^2
			\begin{pmatrix}
				-\Delta & \Omega_0 \\
				\Omega_0 & \Delta
			\end{pmatrix}^2
			\begin{pmatrix}
				\tilde{c}_g(t) \\
				\tilde{c}_e(t)
			\end{pmatrix}
			= \\
			&= -\frac{1}{4}
			\begin{pmatrix}
				\Delta^2 + \Omega_0^2 & 0 \\
				0 & \Delta^2 + \Omega_0^2
			\end{pmatrix}
			\begin{pmatrix}
				\tilde{c}_g(t) \\
				\tilde{c}_e(t)
			\end{pmatrix} \equiv
			-\frac{1}{4}
			\begin{pmatrix}
				\Omega^2 & 0 \\
				0 & \Omega^2
			\end{pmatrix}
			\begin{pmatrix}
				\tilde{c}_g(t) \\
				\tilde{c}_e(t)
			\end{pmatrix},
		\end{align*}
	
	where in the last step, the generalized Rabi oscillation $\Omega$ was used.\\
	\\
	It can be shown, that these equations are solved by
		\begin{align*}
			\tilde{c}_g(t) = -i\frac{\Omega_0}{\Omega} \sin\lr{ \frac{\Omega t}{2} } \tilde{c}_e(0) + \lreck{ \cos\lr{ \frac{\Omega t}{2} } + i \frac{\Delta}{\Omega} \sin\lr{ \frac{\Omega t}{2} } } \tilde{c}_g(0) \\
			\tilde{c}_e(t) = -i\frac{\Omega_0}{\Omega} \sin\lr{ \frac{\Omega t}{2} } \tilde{c}_g(0) + \lreck{ \cos\lr{ \frac{\Omega t}{2} } - i \frac{\Delta}{\Omega} \sin\lr{ \frac{\Omega t}{2} } } \tilde{c}_e(0).
		\end{align*}

	The occupation probabilities can be calculated by plugging in the substitutions $c, \tilde{c}$:
		\begin{align*}
			P_e(t) &:= |C_e(t)|^2 = |c_e(t) e^{-i\omega_e t}|^2 = |c_e(t)|^2 = |\tilde{c}_e e^{i\Delta t/2}|^2 = |\tilde{c}_e(t)|^2 \\
			P_g(t) &:= |C_g(t)|^2 = ... = |\tilde{c}_g(t)|^2
		\end{align*}

	Using the initial conditions $C_g(0)=1$ and $C_e(0)=0$ yields:
		\begin{align*}
			P_e(t) &= |\tilde{c}_e(t)|^2 = \left|-i\frac{\Omega_0}{\Omega} \sin\lr{ \frac{\Omega t}{2} } \right|^2 = \lr{ \frac{\Omega_0}{\Omega} }^2 \sin^2 \lr{ \frac{\Omega t}{2} } \qquad \blacksquare \\
			P_g(t) &= |\tilde{c}_g(t)|^2 = ... = \cos^2\lr{ \frac{\Omega t}{2} } + \lr{ \frac{\Delta}{\Omega} }^2 \sin^2\lr{ \frac{\Omega t}{2} } \qquad \blacksquare
		\end{align*}
\hrule
}


\newpage


\section*{Appendix B: ABCD Formalism}
The ABCD formalism is a formalism to describe light rays and the effect of optical systems on them.\\
\\
The light ray is described by the vector
	\begin{equation}
		\begin{pmatrix}
			x \\ \theta
		\end{pmatrix},
	\end{equation}

where $x$ is the distance to the optical axis and $\theta$ the angle between the light ray and the optical axis.\\
\\
\underline{Convention}: Light rays, that travel away from the optical axis have an angle $\theta > 0$ and those who travel towards the axis $\theta < 0$.\\
Concave mirrors have a negative curvature $R<0$ and convex mirrors a positive one $R>0$.\\
\\
Important matrices are:
	\begin{itemize}
		\item \textbf{Propagation}\\
		The propagation of a light ray in a dielectric medium by the length $d$ is described by the matrix
			\begin{equation}
				\begin{pmatrix}
					1 & d \\
					0 & 1
				\end{pmatrix}.
			\end{equation}


		\item \textbf{Refraction}\\
		The refraction of a light beam that travels in the medium $n_1$ and enters medium $n_2$ is:
			\begin{equation}
				\begin{pmatrix}
					1 & 0 \\
					0 & \frac{n_1}{n_2}
				\end{pmatrix}
			\end{equation}


		\item \textbf{Refraction at a sphere with curvature $R$}	\\
		The refraction at a sphere with curvature $R$ is described by
			\begin{equation}
				\begin{pmatrix}
					1 & 0 \\
					-\frac{n_2 - n_1}{n_2 \cdot R} & \frac{n_1}{n_2}
				\end{pmatrix},
			\end{equation}
		where $R>0$ for a convex and $R<0$ for a concave mirror.


		\item \textbf{Mirror}\\
		The reflection off a mirror is described by
			\begin{equation}
				\begin{pmatrix}
					1 & 0 \\
					\frac{2}{R} & 1
				\end{pmatrix},
			\end{equation}
		
		where $R>0$ for a convex and $R<0$ for a concave mirror.


		\item \textbf{Lens}\\
		A thick lens of width $d$ and curvatures $R_1$ and $R_2$, where $n_1$ is the refractive index of the medium outside the lens and $n_2$ inside the lens, is described by the matrix
			\begin{equation}
				\begin{pmatrix}
					1 - d \frac{n_2-n_1}{n_1}\frac{1}{R_1} & d \frac{n_1}{n_2} \\
					\frac{n_2-n_1}{n_1} \lr{ \frac{1}{R_2} - \frac{1}{R_1} - \frac{n_2-n_1}{n_1}\frac{d}{R_1 R_2} } & 1 + \frac{n_2-n_1}{n_1}\frac{d}{R_2}
				\end{pmatrix}
				\stackrel{\comment{d\to 0}}{\longrightarrow}
				\begin{pmatrix}
					1 & 0 \\
					\frac{n_2-n_1}{n_1} \lr{ \frac{1}{R_2} - \frac{1}{R_1} } & 1
				\end{pmatrix}
			\end{equation}
	\end{itemize}


\newpage


{\color{violet}
\hrule
	\fbox{Derivation}\\
	ges: Propagation matrix\\
	\\
	A light beam inside a dielectric medium propagates without refraction, so
		\begin{align*}
			\theta' &= \theta.
		\end{align*}
	
	Let $x$ be the distance from the optical axis at the beginning of the propagation and $x'$ the distance after propagation.\\
	Consider the triangle spanned by the propagated light beam as well as the line perpendicular and paralel to the optical axis. Formulate the tangent for the incident angle $\theta$:
		\begin{align*}
			\tan\theta &\approx \theta = \frac{x'-x}{d} \\
			\Leftrightarrow x' &= x + d \cdot \theta
		\end{align*}
	
	Reformulate the found equations for $x'(x,\theta)$ and $\theta'(x,\theta)$ as matrix multiplication:
		\begin{align*}
			\begin{pmatrix}
				x' \\ \theta'
			\end{pmatrix}
			=
			\begin{pmatrix}
				1 & d \\
				0 & 1
			\end{pmatrix}
			\cdot
			\begin{pmatrix}
				x \\ \theta
			\end{pmatrix}
			\qquad \blacksquare
		\end{align*}


	\fbox{Derivation}\\
	ges: Refraction\\
	\\
	A light beam, that gets refracted at a medium border, keeps its immediate distance to the optical axis
		\begin{align*}
			x' &= x.
		\end{align*}
	
	The change of angle is described by Snell's law, which in the paraxial approximation simplifies to
		\begin{align*}
			n_1 \sin\theta &= n_2 \sin\theta' \\
			\Leftrightarrow n_1 \theta &= n_2 \theta' \\
			\Leftrightarrow \theta' &= \frac{n_1}{n_2} \cdot \theta.
		\end{align*}

	Reformulate the found equations for $x'(x,\theta)$ and $\theta'(x,\theta)$ as matrix multiplication:
		\begin{align*}
			\begin{pmatrix}
				x' \\ \theta'
			\end{pmatrix}
			=
			\begin{pmatrix}
				1 & 0 \\
				0 & \frac{n_1}{n_2}
			\end{pmatrix}
			\cdot
			\begin{pmatrix}
				x \\ \theta
			\end{pmatrix}
			\qquad \blacksquare
		\end{align*}


	\fbox{Derivation}\\
	ges: Mirror\\
	\\
	A light beam, that gets refracted at a medium border, keeps its immediate distance to the optical axis
		\begin{align*}
			x' &= x.
		\end{align*}
	

\newpage


	Picture a light ray that falls onto a spherical mirror.\\
	Let the angle between the light ray and the optical axis be $\theta$, $\varphi$ the angle of incidence at the mirror, $\alpha$ the angle between the normal and the optical axis and $\theta'$ the angle between the reflected ray and the optical axis.\\
	\\
	Formulate angle equations to connect $\theta'$ with $x$ and $\theta$.\\
	Notice, that the incoming ray moves away from the optical axis, thus $\theta > 0$ and that the outgoing ray moves towards the optical axis, thus $0>\theta' \equiv -\theta'$:
		\begin{align*}
			\alpha &= \varphi - \theta \\
			\theta' &= 2\varphi - \theta
		\end{align*}
	
	Rearranging the first equation for $\varphi$, inserting it into the second one and making use of $\tan\alpha \approx \alpha = \frac{x}{-R}$ yields:
		\begin{align*}
			\theta' &= 2\lr{ -\frac{x}{R} + \theta } - \theta = -\frac{2}{R} \cdot x + 1 \cdot \theta
		\end{align*}

	Reformulate the found equations for $x'(x,\theta)$ and $\theta'(x,\theta)$ as matrix multiplication:
		\begin{align*}
			\begin{pmatrix}
				x' \\ \theta'
			\end{pmatrix}
			=
			\begin{pmatrix}
				1 & 0 \\
				-\frac{2}{R} & 1
			\end{pmatrix}
			\cdot
			\begin{pmatrix}
				x \\ \theta
			\end{pmatrix}
			\qquad \blacksquare
		\end{align*}


	\fbox{Derivation}\\
	zz: Refraction at a sphere with curvature $R$\\
	\\
	Consider a concave spherical border between the medii $n_1$ and $n_2$ with the curvature $R$.\\
	Let the angle between the light ray and the optical axis be $\theta$, $\varphi$ the angle of incidence at the mirror, $\varphi'$ the angle between the normal and the optical axis and $\theta'$ the angle between the reflected ray and the optical axis.\\
	\\
	Snell's law, under paraxial approximation, yields the relation between $\varphi$ and $\varphi'$:
		\begin{align*}
			n_1 \sin \varphi &= n_2 \sin \varphi' \\
			\Leftrightarrow n_1 \cdot \varphi &= n_2 \cdot \varphi' \\
			\Leftrightarrow \varphi' &= \frac{n_1}{n_2} \cdot \varphi
		\end{align*}

	Formulate angle equations to connect $\theta'$ with $x$ and $\theta$:
		\begin{align*}
			\alpha &= \theta' + \varphi' \\
			\alpha &= \theta + \varphi \\
			\alpha &\approx -\frac{x}{R}
		\end{align*}


\newpage


	Rearrange the first equation for $\theta'$ and insert $\alpha, \varphi'$ and $\varphi$:
		\begin{align*}
			\theta' &= \alpha - \varphi' = \\
			&= -\frac{x}{R} - \lreck{ \frac{n_1}{n_2} \cdot \varphi } = \\
			&= -\frac{x}{R} - \lreck{ \frac{n_1}{n_2} \cdot \lr{ \alpha - \theta } } = \\
			&= -\frac{x}{R} - \lreck{ \frac{n_1}{n_2} \cdot \lr{ -\frac{x}{R} - \theta } } = \\
			&= - \frac{x}{R} + \frac{n_1}{n_2} \cdot \frac{x}{R} + \frac{n_1}{n_2} \cdot \theta = \\
			&= -\frac{(n_2 - n_1)}{n_2 \cdot R} \cdot x + \frac{n_1}{n_2} \cdot \theta
		\end{align*}
	
	Reformulate the found equations for $x'(x,\theta)$ and $\theta'(x,\theta)$ as matrix multiplications:
		\begin{align*}
			\begin{pmatrix}
				x' \\ \theta'
			\end{pmatrix}
			&=
			\begin{pmatrix}
				1 & 0 \\
				-\frac{(n_2 - n_1)}{n_2 \cdot R} & \frac{n_1}{n_2}
			\end{pmatrix}
			\cdot
			\begin{pmatrix}
				x \\ \theta
			\end{pmatrix}
			\qquad \blacksquare
		\end{align*}






	\fbox{Derivation}\\
	zz: Thin lens\\
	\\
	The matrix for a thin lens can be described as a diffraction at a sphere with curvature $R_1$, a translation of length $d$ and another diffraction at a sphere with curvature $R_2$. Notice, that the order of matrices is ''reversed'':
		\begin{align*}
			&
			\begin{pmatrix}
				1 & 0 \\
				-\frac{(n_1 - n_2)}{-n_1 R_2} & \frac{n_2}{n_1}
			\end{pmatrix}
			\cdot
			\begin{pmatrix}
				1 & d \\
				0 & 1
			\end{pmatrix}
			\cdot 
			\begin{pmatrix}
				1 & 0 \\
				-\frac{(n_2 - n_1)}{n_2 R_1} & \frac{n_1}{n_2}
			\end{pmatrix} = \\
			&=
			\begin{pmatrix}
				1 & 0 \\
				\frac{(n_1 - n_2)}{n_1 R_2} & \frac{n_1}{n_2}
			\end{pmatrix}
			\cdot
			\begin{pmatrix}
				1 - d \cdot \frac{n_2 - n_1}{n_2 R_1} & d \cdot \frac{n_1}{n_2} \\
				- \frac{n_2 - n_1}{n_2 R_1} & \frac{n_1}{n_2}
			\end{pmatrix} = \\
			&=
			\begin{pmatrix}
				1 - d \cdot \frac{n_2 - n_1}{n_2 R_1} & d \cdot \frac{n_1}{n_2} \\
				\frac{n_1 - n_2}{n_1 R_2} \cdot \lr{ 1 - d \cdot \frac{n_2 - n_1}{n_2 R_1} } - \frac{n_2}{n_1} \cdot \frac{n_2 - n_1}{n_2 R_1} & \frac{n_1 - n_2}{n_1 R_2} \cdot d \cdot \frac{n_1}{n_2} + 1
			\end{pmatrix} \longrightarrow \\
			&\stackrel{\comment{d\to\infty}}{\longrightarrow}
			\begin{pmatrix}
				1 & 0 \\
				\frac{n_1 - n_2}{n_1 R_2} - \frac{n_2}{n_1} \cdot \frac{n_2 - n_1}{n_2 R_1} & 1
			\end{pmatrix} = \\
			&=
			\begin{pmatrix}
				1 & 0 \\
				\lr{ 1 - \frac{n_2}{n_1} } \lr{ \frac{1}{R_1} + \frac{1}{R_2} } & 0
			\end{pmatrix} \equiv \\
			&\equiv
			\begin{pmatrix}
				1 & 0 \\
				-\frac{1}{f} & 0
			\end{pmatrix}
			\qquad \blacksquare
		\end{align*}
\hrule
}


\newpage


\section{Stichwortverzeichnis}
\begin{xltabular}{\linewidth}{p{6cm}|X}
	\hline\hline
		Begriff & Erklärung \\
	\hline
		bare states & Eigenstates $\{\ket{g,n}, \ket{e,n-1}\}$ of the Jaynes-Cummings-Hamiltonian without the interaction term $\hat{H}_0$ \\
	\hline
		cavity dumping & Removal of one of the cavity mirros, which outputs all photons at once \\
	\hline
		coherent regime &  \\
	\hline
		coherence function, temporal & Autocorrelation function of the electromagnetic wave at a fixed point
			$$G(\vec{r},\tau) := \lim\limits_{T\to\infty} \frac{1}{2T} \int_{-T}^{+T} \diff t \, U(t) \cdot U(t+\tau)$$ \\
	\hline
		coherence length & Length $L_\text{coh}$ that relates to the coherence time $\tau_\text{coh}$ via
			$$L_\text{coh} = c_m \cdot \tau_\text{coh}$$ \\
	\hline
		coherence time & Shift $\tau_\text{coh}$ for which the degree of temporal coherence drops to $1/2$:
			$$g(\tau_\text{coh}) = \frac{1}{2}$$ \\
	\hline
		complex spatial coherence & Standardised spatial coherence function \\
	\hline
		complex temporal coherence & Standardised coherence function
			$$\gamma(\vec{r},\tau) := \frac{G(\vec{r},\tau)}{G(\vec{r},0)}$$ \\
	\hline
		cross section & \\
	\hline
		dipole approximation &  \\
	\hline
		degree of spatial coherence & Absolute value of the complex spatial coherence
			$$|\gamma(\vec{r}_1,\vec{r}_2,t)|$$ \\
	\hline
		degree of temporal coherence & Absolute value of the complex temporal coherence
			$$g(\tau) := |\gamma(\vec{r},\tau)|$$ \\
	\hline
		dressed states & Eigenstates of the Jaynes-Cummings-Hamiltonian's subspace $\hat{H}_n$ 
			$$\ket{+,n}, \ket{-,n}$$ \\
	\hline
		Einstein coefficients & Coefficients $A, B_{21}, B_{12}$ that describe the rate of the processes spontaneous emission, stimulated emission and absorption \\
	\hline
		frequency pulling &  \\
	\hline
		gain-switchting &  \\
	\hline
		Glauber-Sudarshan P-function & The P-function
			$$P(\alpha) := e^{|\alpha|^2 / 2} \int \diff^2 u \, e^{|u|^2} \braket{ -u | \hat{\rho} | u } e^{u \alpha^\star - u \alpha^\star}$$ 
		serves as a classification for non-classical light.\\
		& If 
			$$P<0$$,
		the state $\hat{\rho}$ has no classical equivalent \\
	\hline
		Gouy-phase & Term $\zeta(z)$ in the phase of a Gaussian beam \\
	\hline
		Hanburry-Brown-Twiss Interferometer &  \\
	\hline
		Helmholtz equation & Wave equation under the assumption of a monochromatic wave $U(\vec{r},t) = U(\vec{r}) \cdot e^{i\omega t}$:
			$$\lr{ \vec{\nabla}^2 + k^2 } U(r) = 0$$ \\
	\hline
		Helmholtz equation, paraxial & Helmholtz equation under the assumptions \\
		& - a field near and almost parallel to the optical axis $z$ \\
		& - slowly varying envelope approximation $\partial_z A$
			$$\lr{ \vec{\nabla}_x^2 + \vec{\nabla}_y^2 -2ik \partial_z } A(\vec{r}) = 0$$ \\
	\hline
		hole burning &  \\
	\hline
		incoherent regime &  \\
	\hline
		inhomogeneous, paraxial wave equation & eq. that describes the E-field in the classical description of the light-matter interaction under paraxial and slowly varying E-field approximation \\
	\hline
		Kerr lensing &  \\
	\hline
		laser medium & Linear, isotropic, non-magnetic medium which can contain dopants (e.g. atoms) for amplification \\
	\hline
		length gauge & Gauge, which keeps the vector potential in the quantum mechanical light matter interaction Hamiltonian \\
	\hline
		line shape function &  \\
	\hline
		line width &  \\
	\hline
		Lorenz function & Factor in the absorption coefficent in the semi-classical description of light-matter interaction
			$$L(\omega) := \frac{1/\tau_P}{ (\omega - \omega_0)^2 + 1/\tau_P^2 }$$ \\
	\hline
		Mandel Q parameter & Measure for the deviation of the photon statistic of a state from the Poisson distribution
			$$Q := \frac{(\Delta n)^2}{\braket{\hat{n}}} - 1$$ \\
			& - If $Q<0$, one calls the state as \textbf{sub-Poissonian} \\
			& - If $Q=0$, one calls the state \textbf{poissonian} \\
			& - If $Q>0$, one calls the state \textbf{super-poissonian} and classical \\
	\hline
		master equation &  \\
	\hline
		master equation, optical &  \\
	\hline
		mode locking &  \\
	\hline
		mode pulling & Shift of the lasing frequency from the resonator frequency towards the atomic transition frequency \\
	\hline
		mode selection & Process of isolating a specific mode from multi-mode lasing through e.g.\\
		& - grating \\
		& - Littrow arrangement (prism + mirror) \\
		& - Fabry-Pérot etalon \\
		& - waveplate + polarisers \\
	\hline
		multi-mode lasing &  \\
	\hline
		optical oscillator & Amplifier combined with a positive feedback \\
		& A laser is an optical oscillator if the threshold and phase condition are met \\
	\hline
		oscillator strength &  \\
	\hline
		population inversion & Difference between the population densities of the two-state atom
			$$\Delta n := n_2 - n_1$$
		where $n_i := n_0 \cdot |c_i|^2$ with the density of atoms $n_0$ and the probability for state $i$ $|c_i|^2$\\
		& It appears in the gain coefficient of a laser\\
	\hline
		Q-switching & cavity modulation which reduces the losses in the laser medium for a short time, thus creating pulsed light \\
	\hline
		quadrature operators & Linear combinations of the annhilation and creation operator that are proportional to the position/momentum operator
			{
			\begin{align*}
				\hat{X}_1 &:= \frac{\hat{a} + \hat{a}^\dagger}{2} \propto \hat{q} \\
				\hat{X}_2 &:= \frac{\hat{a} - \hat{a}^\dagger}{2i} \propto \hat{p}
			\end{align*}
			}
			\\
	\hline
		resonator, cold & Resonator with a population inversion far below the threshold
			$$\Delta n \ll \Delta n_\text{thr}$$ \\
	\hline
		resonator, hot & Resonator with a population inversion far above the threshold
			$$\Delta n \gg \Delta n_\text{thr}$$ \\
	\hline
		rotating wave approximation (RWA) &  \\
	\hline
		saturation intensity & \\
	\hline
		Shawlow-Townes Limit & Expression for the laser linedwidth above the threshold condition
			$$\Delta\nu_L \propto \frac{1}{P}$$
		with the pump power $P$ \\
	\hline
		single-mode lasing & \\
	\hline
		stability diagramm &  \\
	\hline
		threshold condition & Population inversion that has to be exceeded for the gain to be greater than the losses and thus laser oscillations to occur
			$$\Delta n_\text{thr} = \frac{\alpha}{\sigma(\nu)}$$ \\
	\hline
		total gain & The total gain is the relation between the intensity after propagation and the initial intensity
			$$G := \frac{I(z=d)}{I(z=0)}$$ \\
	\hline
		velocity gauge & Gauge, which keeps the E-field term \\
	\hline
		wave front & Surface of constant phase \\
	\hline
		Wiener-Khinchin theorem & The power spectrum of a light field $W(\nu)$ is the Fourier Transform of the temporal coherence $\Gamma(\tau)$
			$$W(\nu) = \int_{-\infty}^\infty \diff\tau \, \Gamma(\tau) \cdot e^{i\nu\tau}$$ \\
	\hline
		visibility & Property of an interference pattern
			$$V = \frac{ I_\text{max} - I_\text{min} }{ I_\text{max} + I_\text{min} }$$ \\
	\hline\hline
\end{xltabular}
\end{document}
